\section{Proposed Protocol Extensions}
\label{sec:proposals}
\setlength{\emergencystretch}{3em}

This section presents the three protocol extensions the
survey introduces as Contributions C5, C6, and C7
(\S\ref{sec:intro:contributions}). The proposals are not
the survey's primary contribution; the rubric and the
comparison artefacts are. They do fall out directly
from the open-problems agenda of \S\ref{sec:openproblems},
so we include them as derived contributions. That is more
honest than leaving them in the closing-paragraph
speculation that typically ends a survey of this kind.

The three proposals are co-designed: Contribution C5 (the
OWA naming scheme, \S\ref{sec:proposals:owa}) provides the
identifier substrate; Contribution C6 (the AGRP guardrails
envelope, \S\ref{sec:proposals:agrp}) provides the
span-level egress-attestation envelope; Contribution C7
(the ACSP context selection envelope,
\S\ref{sec:proposals:acsp}) provides the
item-level selection envelope that decides which items
cross the boundary before AGRP cleans the spans within
them. The boundary-detection primitive that both AGRP
(\S\ref{sec:proposals:agrp:boundary}) and ACSP rely on is
derived directly from prefix comparison of OWA URIs. The
closing
subsection~(\S\ref{sec:proposals:codesign}) explains why
the co-design is intrinsic rather than incidental, and
why no pair of the three proposals is sufficient without
the third.

All three proposals are at specification-sketch stage at
the 2026-Q1 cutoff, with no validated protocol
implementation in the cutoff corpus. Because that phrase is
easy to read past, we state exactly what it excludes. There
is no formal proof of any security property claimed here.
Nothing proves that AGRP's tier semantics are unforgeable.
Nothing proves that ACSP's verification steps are sound or
complete against a stated adversary. Nothing proves that
OWA's prefix comparison decides the boundary correctly under
adversarial naming. There are no measured overheads either:
no latency distribution, no throughput figure, and no size
cost for either envelope on a real workload. And there is one partial
reference implementation, not three. A post-cutoff PACMS
artefact (\S\ref{sec:proposals:acsp:pacms}) instantiates the
selector side of part of ACSP for single-organisation use. It
carries none of its signatures, receiver verification, or
inter-organisational protocol. OWA and AGRP have no
implementation at all.

What the proposals do carry is precision: an ABNF for OWA,
schemas and normative verification steps for both envelopes,
and a stated failure mode for each. Each proposal's open
questions subsection names the ones we consider
load-bearing. Two of them are properties a deployer would
want to know before adopting. ACSP's selection log discloses
candidate-set membership across a boundary
(\S\ref{sec:proposals:acsp:open}, Q3). Every AGRP guarantee
is conditional on a content classifier that will be wrong in
both directions (\S\ref{sec:proposals:agrp:open}, Q4). Both
are analysed rather than listed.

All three are candidates for the next AAIF and W3C revision
cycles that the governance analysis of
\S\ref{sec:governance:survival} anticipated. The survey is
honest about this maturity gap (R3 in our source-tier
discipline, \S\ref{sec:intro:notwhat}): the proposals are
reference designs, not validated specifications, and the
section positions them accordingly.

The security-gating code and deterministic feasibility
harness described later were produced in a post-cutoff
collaboration artefact (August 2026). They are reported to
make the proposed evaluation reproducible, but they do not
change the 2026-Q1 literature corpus, differentiation
tables, or rubric scores. No count from that harness is a
survey result.

The contemporary literature on inter-agent naming,
guardrails, and context selection is more developed than
the early-2025 baseline of the prior surveys; the three
differentiation tables of this section
(Table~\ref{tab:owa-differentiation},
Table~\ref{tab:agrp-differentiation}, and
Table~\ref{tab:acsp-differentiation}) are the centrepiece
of the proposals' positioning, and the reader who is only
checking that the proposals are non-trivial against the
literature should turn directly to those tables.

\subsection{OWA Scoping Scheme (Contribution C5)}
\label{sec:proposals:owa}

\subsubsection{Problem statement}
\label{sec:proposals:owa:problem}

The contemporary inter-agent protocol landscape has
accumulated multiple naming and addressing schemes, none of
which carve out an explicit \emph{workspace} tier between
the organizational trust root and the agent identifier. A2A
identifies agents by an AgentCard URI hosted at a
well-known path of the agent's host; ANP uses DIDs; the
Agent Identity URI Scheme~\cite{agent_uri_scheme} introduces
a stable \texttt{agent://} URI but leaves workspace
tenancy unaddressed; the Agent Name Service IETF
draft~\cite{ans_ietf} flattens organizational identity into
a single \texttt{Provider} field that conflates the
organization with the workspace.

The problem this collapses is the empirical reality that
enterprises run \emph{multiple agents performing
overlapping capabilities under different workspace
policies}: a billing agent under the finance workspace, a
billing agent under the customer-support workspace, and a
billing agent under the engineering workspace are three
distinct agents that share a capability descriptor but
that exist under three different governance regimes.
The regimes differ in access scopes, audit policies,
escalation paths, and egress filters.
Existing schemes either collapse the three into a single
identifier (losing the workspace distinction) or
hostname-mangle the workspace into the FQDN
(\texttt{finance.acme.com}). The subdomain form works for
a single deployment but has three specific drawbacks that
we develop in \S\ref{sec:proposals:owa:related}:
it fragments the organizational trust root across
subdomains, it couples workspace lifecycle to DNS-zone
administration, and it does not accommodate the
cardinality and churn rates at which workspaces are
created and retired in production.

The trust-bootstrapping problem (P7 of
\S\ref{sec:openproblems:trust}) overlays this naming gap:
agents from heterogeneous registries that need to
cross-trust each other have no normative way to express
which workspace the trust applies to. A2A AgentCards
under one organization's well-known endpoint and ANP DIDs
under another organization's DID resolver can not name
the same workspace without an out-of-band convention.

\subsubsection{Related work and differentiation}
\label{sec:proposals:owa:related}

Table~\ref{tab:owa-differentiation} compares OWA against
the eight most-relevant prior works in the inter-agent
naming space. The comparison covers two dimensions of
naming-scheme design (does the scheme have an
organizational tier? a workspace tier?), two dimensions
of semantic content (does the scheme carry capability
semantics? does it separate identity from location?), and
the co-design relationship with the survey's other
proposal.

\begin{table*}[t]
  \centering
  \caption{OWA differentiation matrix. The five-property
    argument: no prior scheme simultaneously provides all
    five of (a) DNS-anchored org trust root, (b) first-class
    workspace tier, (c) capability semantics, (d) stable
    identity decoupled from location, and (e) co-design with
    cross-boundary guardrails.}
  \label{tab:owa-differentiation}
  \footnotesize
  \resizebox{\textwidth}{!}{%
  \begin{tabular}{p{3.2cm}p{2.4cm}cccc p{3.0cm}}
    \toprule
    Prior work & Form & Org? & Workspace? & Capability? & Identity~$\neq$~location? & OWA improvement \\
    \midrule
    \texttt{agent://} URI~\cite{agent_uri_scheme}
      & \texttt{agent://root/cap/id}
      & yes & \textbf{no} & yes & yes
      & adds workspace tier under distinct \texttt{agent+owa://} scheme suffix; reuses trust-root and capability-path layers \\
    ANS IETF draft~\cite{ans_ietf}
      & \texttt{proto://id.cap.provider.vN}
      & yes (Provider) & \textbf{no (flat)} & yes & partial
      & separates org tenancy from workspace tenancy that ANS collapses \\
    A2A AgentCard~\cite{spec_a2a}
      & \texttt{https://host/.well-known/agent.json}
      & implicit & no & no & \textbf{no (coupled)}
      & stable identity decoupled from host (Saltzer 1982) \\
    W3C DID
      & \texttt{did:method:id}
      & no & no & \textbf{no} & yes
      & adds capability and tenancy that DIDs leave to higher layers \\
    Email~(RFC~5321, real SMTP)
      & \texttt{local@domain}
      & yes & only via \texttt{+suffix} hack & \textbf{no} & partial (MX)
      & full email transport has prohibitive latency / metadata / spam dynamics for inter-agent \\
    Email-syntax IAM identity~\cite{gcp_service_accounts, entra_service_principals}
      & \texttt{name@project.iam.gsa.com} (GCP) / \texttt{name@tenant.onmicrosoft.com} (Entra)
      & yes (DNS) & via project/tenant, no native workspace tier & \textbf{no} & yes
      & cross-vendor protocol-level standardisation; capability path; explicit workspace tier \\
    HTTP URI
      & \texttt{https://host/path}
      & implicit & path-only, no semantics & path-only & \textbf{no}
      & separates identity from location \\
    AWCP~\cite{awcp}
      & delegation protocol
      & n/a & as behaviour & partial & n/a
      & names workspaces; AWCP delegates within them (complementary) \\
    Ruflo federation~\cite{ruflo}
      & mTLS + ed25519 + WS
      & yes (per-machine) & \textbf{no normative scheme} & no & partial
      & gives Ruflo's federation a normative workspace name \\
    AI Catalog~\cite{ai_catalog}
      & \texttt{urn:air:pub:ns:name}
      & yes (publisher) & namespace tier (optional, hierarchical) & \textbf{no} & yes
      & adds ABNF, canonical form, prefix-matchability; OWA workspace is a normative tier, not optional \\
    AgentDNS~\cite{agentdns}
      & \texttt{agentdns://org/cat/name}
      & yes (registered at a root) & \textbf{no} & yes (category) & yes (resolved metadata)
      & DNS-bootstrapped org root instead of a registering authority; workspace tier; no proxy in the invocation path \\
    AI Identity gaps survey~\cite{ai_identity_gaps}
      & survey
      & flags the gap & flags the gap & flags the gap & flags the gap
      & OWA is the proposal the survey calls for \\
    \bottomrule
  \end{tabular}%
  }
\end{table*}

Three cells of the table warrant prose elaboration because
they carry the differentiation argument most heavily.

\paragraph{The agent:// scheme is the closest precedent.}
Rodriguez~\cite{agent_uri_scheme} introduced
\texttt{agent://} explicitly to separate identity from
location, citing Saltzer's 1982 analysis that names should
be resolved to addresses rather than used as addresses. The
scheme's three-tier structure
(\texttt{trust-root/}\allowbreak
\texttt{capability-path/}\allowbreak\texttt{agent-id}) covers
organizational identity and capability semantics with the
stable-identity discipline that A2A's AgentCard URI lacks.
OWA adopts the scheme's identity discipline verbatim, adopts
its trust-root layer (DNS-anchored organizational identity)
and its capability-path layer (hierarchical, prefix-matched),
and inserts a workspace tier between the trust root and the
capability path.

OWA is not a transparent extension of \texttt{agent://} at
the URI-parsing layer. An unmodified \texttt{agent://}
resolver that receives an OWA URI does not error cleanly:
the reference parser treats the trailing typed identifier
as the agent-id and folds the interstitial
\texttt{/w/<workspace>/c/<cap>/a/} into the capability
path, yielding a DHT key over the wrong canonical name and
a resolver query that will never hit the intended record.
The single-character delimiters \texttt{/w/}, \texttt{/c/},
and \texttt{/a/} also collide with legal single-character
capability segments in
\texttt{agent://}'s published test vectors, so a parser
cannot distinguish an OWA delimiter from a capability
segment without explicit knowledge of the OWA grammar.
Because of both, we introduce OWA as a distinct URI scheme
suffix, \texttt{agent+owa://}, following the
\texttt{git+ssh://} and \texttt{stratum+tcp://} convention:
the suffix signals to any parser that OWA-aware path
grammar is required, and \texttt{agent://}-only parsers
reject the URI cleanly rather than silently mis-parsing.
Backwards compatibility at the \emph{resolver} layer is
preserved: an organization's \texttt{agent://} DNS root,
its \texttt{.well-known/agent-keys.json}, and its
attestation infrastructure serve both schemes unchanged.
Only URI parsers and DHT-key-deriving components need to
be OWA-aware.

\paragraph{ANS conflates the workspace into Provider.}
The ANS IETF draft~\cite{ans_ietf} extends DNS conventions
with capability-aware resolution, producing identifiers
like \texttt{a2a://textProcessor.}\allowbreak
\texttt{DocumentTranslation.}\allowbreak
\texttt{AcmeCorp.v2.1.hipaa}.
The \texttt{Provider} field carries the organizational
identity but flattens any intra-organizational tenancy
into the same flat field. An enterprise with three billing
agents under three workspaces must either invent provider
sub-names (\texttt{AcmeCorp-Finance},
\texttt{AcmeCorp-Support}, \texttt{AcmeCorp-Engineering})
or rely on the \texttt{AgentID} field to disambiguate.
Provider sub-names couple workspace identity to provider
naming, and they make workspace renaming expensive.
Overloading \texttt{AgentID} loses the
workspace-as-policy-boundary that the next subsection's
AGRP proposal requires. OWA gives the workspace its own
URI segment and resolves both problems.

\paragraph{Ruflo demonstrates the federation-without-naming gap.}
Ruflo~\cite{ruflo} ships cross-machine agent federation
with mTLS+ed25519 transport, an optional WireGuard mesh,
and a 14-type PII detection pipeline applied per
trust-level (BLOCK / REDACT / HASH / PASS). The federation
works at deployment time, but the documentation does not
specify a normative agent name that travels with the
federated identity. Trust level is computed by a
trust-scoring engine inside one orchestrator rather than
expressed in the URI. OWA is the missing piece that would
let Ruflo-style federation compose across orchestrators
who don't share a trust-scoring runtime: an OWA URI
carries the organizational and workspace identity into the
federation envelope, and the receiver can apply
trust policies keyed on the URI's components.

\paragraph{AI Catalog adopts the identity/metadata split independently.}
The AI Catalog working group~\cite{ai_catalog} (Microsoft,
Google, Cisco, Amazon, Pulse MCP) is developing ongoing
conventions for identifying and discovering AI artifacts.
The scope is broader than agents: a catalog entry may carry
an A2A agent card, an MCP server, a plugin, a dataset, a
model card, or a nested catalog. Agents are one artifact
type among several, and the work does not define a universal
runtime identity for an agent. ADR-0015 recommends
\texttt{urn:air:\{publisher\}:\{namespace\}:\{name\}} for
open or federated deployments, where publisher is a DNS
domain and namespace is optional hierarchical
colon-separated segments. The \texttt{identifier} field
itself stays an open text format, so the URN is a strong
recommendation rather than a settled requirement. The stated
design rationale adopts the identity/metadata split this
survey argues for in \S\ref{sec:proposals:owa:tradeoff}: the
logical name used for discovery and routing is separated
from the cryptographic identity used for trust verification.
NANDA's Migration~010 aligns its \texttt{trust\_manifest}
field to the AI Catalog \texttt{CatalogEntry} schema,
making AI Catalog upstream of NANDA at the
trust-metadata layer (though not yet at the identifier
layer; see \S\ref{sec:proposals:owa:nanda} for the
narrower claim).

AI Catalog's namespace tier (e.g.,
\texttt{urn:air:acme.com:finance:billing-agent}) is a
third data point in the workspace-in-the-URI question: it
places a hierarchical scope inside the name while keeping
tenancy out of the identifier grammar. The tier is closer
to OWA's workspace than to \texttt{agent://}'s flat
capability path, but differs in three respects that keep
OWA's contribution non-redundant. First, the identifier is
deliberately an open text format with no normative ABNF, no
canonical form, and no stated prefix-matchability, which is
the same underspecified path grammar that produced the
silent mis-parse reproduced against OWA in
\S\ref{sec:proposals:owa:related}. Second, the namespace is
optional and carries no workspace-as-policy-boundary
semantics: nothing in the standard ties namespace membership
to egress policy, audit retention, or delegation scope.
Third, the standard's own open issues (issue~59,
proposing that the artifact-type segment move ahead of the
publisher; issue~70, arguing that the trust manifest's MUST
rules cannot be satisfied by open identifier formats)
indicate that the grammar is still under development rather
than settled prior art. We cite it as ongoing work and pin
the citation to a commit for that reason. OWA's
contribution relative to AI Catalog is the normative ABNF,
the canonical form, the prefix-matchable
workspace-as-policy-boundary, and the co-design with AGRP's
boundary detection. The two are not competing for the same
job: AI Catalog is a discovery and trust container across
artifact types, and OWA is a naming discipline for one of
them.

\paragraph{AgentDNS is the centralized answer to the same
question, and the contrast is the trust root.}
AgentDNS~\cite{agentdns} proposes
\texttt{agentdns://org/category/name} as a global service
identifier, resolved by a root server that also runs
discovery, authentication, and settlement. The proposal is
useful evidence for this survey's premise rather than
against it. A second independent group reached the same
conclusion: inter-agent naming needs an organizational
tier, hierarchical semantic categories, and identifiers that
survive a change of physical address. The
\texttt{category} segment is the same move OWA's
\texttt{<cap>} segment makes, and AgentDNS's natural-language
service search is the semantic-discovery appetite
\S\ref{sec:semantic} documents from the other direction.

The designs part company on where trust is rooted, and the
difference has three consequences. First, AgentDNS's
\texttt{org} segment is a name registered with and verified
by the AgentDNS root. OWA's \texttt{<org>} segment is a DNS
domain. Organizational authority is therefore whatever
authority the organization already has over its own domain,
and no new registrar has to exist or stay solvent. Second,
AgentDNS places a service proxy in the invocation path. A
vendor registers a service, the root creates a proxy for it,
and calls route through the root. That makes the root a
mandatory intermediary for traffic as well as for
resolution, and gives it visibility over every inter-agent
call in the ecosystem. OWA resolves names and then gets out
of the way. The message goes directly between the agents,
which is what lets AGRP and ACSP attest at the sender's
boundary rather than at a shared middlebox. Third, AgentDNS
centralizes authentication deliberately, with one
time-bound token accepted by every registered service. That
is a genuine ergonomic gain over per-vendor API keys. It is
also a single point of compromise with the blast radius
of the whole registry, so the registry-poisoning class of
\S\ref{sec:security} lands on it with full force.

Neither position is strictly better, and the honest summary
is that they optimize different things. AgentDNS optimizes
onboarding: one registration, one credential, one bill.
OWA optimizes independence and verifiability. It adds no new
authority and no proxy in the data path. It also carries a
policy boundary a receiver can check from the name alone,
namely the workspace tier, which AgentDNS does not
have. For deployments in
the regulated verticals of \S\ref{sec:adoption:healthcare}
to \ref{sec:adoption:finance}, the proxy and the shared
credential are the two properties most likely to be
disqualifying. That is the case OWA is built for. The
appetite AgentDNS demonstrates is real either way, and a
naming layer that ignores it is unlikely to be adopted.

\paragraph{OWA prefers the path form over the subdomain
form for three specific reasons.}
The most common operator instinct when adding workspace
tenancy to an existing agent identifier is to encode the
workspace as a DNS subdomain of the organizational root,
producing hosts of the form \texttt{finance.acme.com}. The
form works for a single deployment, and we are not aware
of a published inter-agent naming scheme that adopts it,
but it is the alternative that reviewers in early drafts
of this survey most frequently raised, and it warrants
explicit treatment.

\emph{First, trust-atom singularity.} In
\texttt{agent://}~\cite{agent_uri_scheme}, the
organizational trust root runs the key infrastructure that
resolvers, agents, and callers depend on: a single
\texttt{.well-known/agent-keys.json}, a single attestation
issuer, and a single set of signing keys under DNSSEC.
Subdomaining fragments this across N hosts: does
\texttt{finance.acme.com} publish its own
\texttt{.well-known/agent-keys.json}, or inherit
\texttt{acme.com}'s, and how do federated callers verify
which one is authoritative? The path form keeps one
organization as one trust and key atom, with the workspace
as a policy sub-scope beneath it; the trust-root layer of
\texttt{agent://} composes into OWA unchanged.

\emph{Second, DNS zone control.} Creating or renaming a
subdomain is a privileged operation on the organization's
DNS zone: it requires zone-file edits, propagation
timing, and per-zone signing if DNSSEC is enabled. In
enterprises that centralize DNS administration, workspace
lifecycle becomes gated on a ticket queue outside the
workspace owner's control. The path form encodes workspace
as a self-service registry entry that the workspace owner
manages, decoupling workspace lifecycle from DNS-zone
administration.

\emph{Third, workspace cardinality and churn.} Empirically,
workspaces spin up per project, per customer, or per
environment; production deployments carry hundreds to
thousands of workspaces that churn on weekly to monthly
timescales. DNS was engineered for hostnames that change
on the timescale of infrastructure moves, with propagation
delays, per-zone signing costs, and TTL considerations
that assume slow-changing state. A path-segment registry
entry absorbs the cardinality and churn rate that DNS was
not built for.

The three reasons together, not any one of them, justify
the path form over the subdomain form. We do not claim
that OWA composes across organizations more cleanly than
subdomains: neither form solves that problem, and it
remains an open question in
\S\ref{sec:proposals:owa:open}. Nor do we lean on
workspace renaming as a differentiator; CNAME records give
subdomains a mature aliasing mechanism, and the OWA
migration path (\S\ref{sec:proposals:owa:migration}) uses
a workspace-catalogue alias that we do not claim is more
mature than CNAME. The three reasons above are the
substantive ones.

\paragraph{Email-syntax identifiers are the production-validated
precedent, not the \texttt{+suffix} hack.}
The first instinct when reading the table's Email row is
that real RFC 5321~\cite{rfc5321} mail is what the row
dismisses. That is
not the design alternative that warrants serious treatment.
The substantive alternative is the email-syntax identifier
used as an IAM anchor without SMTP transport behind it.
This is the pattern that Google Cloud Service Accounts have
deployed at scale since circa 2014~\cite{gcp_service_accounts}
and that Microsoft Entra mirrors with UPN-shaped
identifiers for service principals and application
identities~\cite{entra_service_principals}. A GCP Service
Account address such as \texttt{billing-agent@finance.iam.gserviceaccount.com}
is identifier syntax; no mailbox accepts mail at it. The
syntactic form is reused for familiarity, DNS-anchored trust,
and IAM-operator habits, but the transport is JWT-signed
HTTPS, not SMTP. The pattern works: it governs the
authorisation surface of cloud workloads for tens of
thousands of organisations on each cloud, has accumulated
ten-plus years of operational hardening, and is the closest
existing analog to what an inter-agent identity standard
might look like.

OWA prefers the path form over the email-syntax form for
three specific reasons. \emph{First}, the email
local-part is two-tier and flat: capability paths such
as \texttt{c/workflow/approval/invoice} cannot be encoded
without inventing a delimiter convention that breaks
standard email parsers, whereas the OWA URI carries
capability hierarchies natively. \emph{Second}, workspace
creation in the email-syntax form requires a DNS
subdomain. That is operationally heavier than a URI path
segment in organisations that do not delegate DNS-zone
control to individual workspace owners. \emph{Third}, an
email-shaped identifier that happens also to be a
deliverable mailbox creates a real conflation risk:
humans who type \texttt{billing-agent@finance.acme.com}
into a mail client either reach a real mailbox (which
must then defend against the dominant phishing surface)
or bounce (UX confusion). The path form
\texttt{agent://\ldots} is unambiguous about not being
mail, which is a deployment-discipline property the
email-shape sacrifices.

The three reasons do not eliminate the email-syntax form's
real advantages. Those advantages are agent-human
integration, trigger-by-email patterns, and existing
IAM-operator familiarity. OWA accommodates them
through optional aliases discussed in
\S\ref{sec:proposals:owa:operational}. The reading we
recommend is that the path form is the canonical
encoding and the email-shaped form is a defined alias for
adoption-friction reasons, not that one form is right and
the other wrong.

\subsubsection{The OWA URI: syntax and component semantics}
\label{sec:proposals:owa:syntax}

The OWA URI conforms to RFC 3986~\cite{rfc3986} and reuses the
resolution and trust infrastructure of the
\texttt{agent://} scheme of~\cite{agent_uri_scheme} under
a distinct scheme suffix, \texttt{agent+owa://}, that
signals OWA-aware path grammar to parsers. The ABNF is:

\begin{verbatim}
owa-uri    = "agent+owa://" org "/w/" workspace "/c/" cap "/a/" agent-id
org        = host                  ; RFC 3986 host (DNS-anchored)
workspace  = 1*( ALPHA / DIGIT / "-" / "_" )
cap        = segment *( "/" segment )
segment    = 1*( ALPHA / DIGIT / "-" / "_" )
agent-id   = "agent_" 26( base32 )  ; TypeID over UUIDv7
\end{verbatim}

A concrete example: \\
\texttt{agent+owa://acme.com/w/finance/c/workflow/approval/}\\
\texttt{invoice/a/agent\_01h455vb4pex5vsknk084sn02q}

The scheme suffix, rather than reusing \texttt{agent://},
is load-bearing. As discussed in
\S\ref{sec:proposals:owa:related}, an unmodified
\texttt{agent://} parser silently mis-parses an OWA URI
into a wrong canonical name, and the \texttt{/w/},
\texttt{/c/}, \texttt{/a/} delimiters collide with legal
single-character capability segments in
\texttt{agent://}'s published test vectors. The
\texttt{agent+owa://} suffix eliminates both risks:
non-aware parsers reject the URI cleanly, and aware
parsers know they are on the OWA grammar. Backwards
compatibility is preserved at the layers where it matters:
the trust-root DNSSEC and key-publication infrastructure,
the capability-path semantics, and the resolution machinery
below the URI parser serve both schemes without change.

The four components carry the following semantics, each
specified normatively rather than left to convention:

\paragraph{\texttt{<org>}: organizational trust root.}
A DNS-anchored hostname identifying the issuing
organization. Resolution authority over the URI follows
DNSSEC for the hostname plus the organization's
chosen agent-resolution mechanism (registry, DHT,
on-chain index). Cross-organizational trust uses the
existing DNS+TLS infrastructure as the trust anchor for
the OWA component.

\paragraph{\texttt{<workspace>}: intra-organizational tenancy boundary.}
A workspace name within the organization. The workspace
is the primary unit of policy enforcement: per-workspace
AGRP egress policies (\S\ref{sec:proposals:agrp}),
per-workspace audit retention, per-workspace
delegation roots. Workspace names are
free-form within the regex above but the organization
SHOULD publish a workspace catalogue at
\texttt{https://<org>/.well-known/workspaces.json} for
discovery.

\paragraph{\texttt{<cap>}: hierarchical capability descriptor.}
Identical semantics to \texttt{agent://} at the capability-path
layer: a hierarchical capability path supporting
prefix-match discovery. The \texttt{agent://} paper's
empirical validation across 369 production tools applies
to the capability-path layer unchanged
\cite{agent_uri_scheme}. The OWA-specific delimiters
\texttt{/w/}, \texttt{/c/}, \texttt{/a/} are consumed by
OWA-aware parsers before the capability-path portion is
extracted, so single-character capability segments remain
legal and are not shadowed by the delimiters.

\paragraph{\texttt{<agent-id>}: stable typed identifier.}
A TypeID-prefixed UUIDv7. Stable across the agent's
lifetime; resolved to a current transport endpoint via
the organization's registry / DHT. Identity is decoupled
from location: an agent that migrates from
\texttt{host-1.acme.com} to \texttt{host-2.acme.com}
retains its OWA URI.

\subsubsection{Resolution, discovery, and federation}
\label{sec:proposals:owa:resolution}

Resolution of an OWA URI to a current transport endpoint
follows the \texttt{agent://} DHT-key pattern with a
specific asymmetry: the DHT-key derivation
(SHA256 over a canonical name) is defined over the org
and agent-id components only, and does \emph{not} fold in
the \texttt{<workspace>} segment. A resolver that receives
the URI

\begin{quote}
\texttt{agent+owa://acme.com/w/finance/c/workflow/approval/}\\
\texttt{invoice/a/agent\_01h455vb4pex5vsknk084sn02q}
\end{quote}

first derives the DHT key from the identity-defining
components, retrieves the endpoint record, and then filters
or verifies the record against the workspace label
\texttt{finance} as an additional matching parameter
carried in the resolver query. The workspace check is a
policy-metadata comparison on the retrieved record, not
part of DHT-key derivation, and the endpoint record itself
carries the workspace attestation. A calling agent
verifies that the record's workspace label matches its
expectation; a receiver that obtains an endpoint record
for the wrong workspace MUST fail closed. This split
between the identity-defining components (in the DHT
key) and the workspace as resolution-time policy metadata
(in the record) is deliberate; we develop the rationale
in \S\ref{sec:proposals:owa:tradeoff}.

Discovery proceeds in three modes:

\begin{description}
  \item[Within-workspace.] A direct lookup against the
    organization's registry filtered by workspace; returns
    the agent population of that workspace.
  \item[Cross-workspace within an org.] A federated lookup
    over the organization's workspace catalogue; the
    catalogue lists workspaces and their resolver
    endpoints. Cross-workspace calls are subject to AGRP
    egress policies (\S\ref{sec:proposals:agrp}).
  \item[Cross-org.] Resolution against the target
    organization's resolver, with the OWA URI carrying
    the full org and workspace identity for policy
    enforcement at both endpoints.
\end{description}

Federation across heterogeneous registries (P7 of
\S\ref{sec:openproblems:trust}) is bootstrapped by the
DNS trust root: an OWA URI carries the
organizational FQDN as its org component, and the
receiver's resolver bootstraps trust through the existing
DNSSEC infrastructure rather than requiring a
cross-registry attestation protocol.

\subsubsection{Why workspace belongs in the URI: a trade-off analysis}
\label{sec:proposals:owa:tradeoff}

A design question worth naming, rather than treating as
settled, is whether the workspace tier belongs in the URI
at all. The \texttt{agent://} discipline is to keep in the
identifier only what is stable and identity-defining, and
to push mutable operational and governance facts (the
transport endpoint, capability metadata,
attestation state) out to the resolved record. By that
same discipline, workspace membership is a candidate for
the resolved record rather than the URI: agents move
between workspaces, workspaces rename, an agent may be a
member of two workspaces, and workspace access lists
change on operational timescales. All three properties are
acknowledged in \S\ref{sec:proposals:owa:open}, and
\S\ref{sec:proposals:owa:resolution} already treats the
workspace label as a matching parameter in the resolver
query rather than as a component of the DHT-key
derivation. Behaviorally, workspace already looks like
resolution-time metadata.

OWA nevertheless carries the workspace in the URI. We
justify this on two grounds, and only two.

\emph{Trust-atom singularity.} The organizational trust
root is a single atom in \texttt{agent://}: one
\texttt{.well-known/}\allowbreak\texttt{agent-keys.json}, one attestation
issuer, one signing-key set under DNSSEC. Keeping the
organization in the URI as the trust root, with workspace
as a path-segment sub-scope beneath it, preserves
that atom. Alternatives that move the workspace out of the
URI must either replicate the trust atom per workspace
(fragmenting keys) or drop workspace scope from the trust
layer entirely (losing the ability to express
workspace-scoped attestations). We are unwilling to pay
either price, and note this is the same singularity that
made the subdomain form unattractive in
\S\ref{sec:proposals:owa:related}.

\emph{Prefix-detection convenience for AGRP.} AGRP
(\S\ref{sec:proposals:agrp}) detects the workspace
boundary of a call by prefix comparison on the sender and
receiver URIs alone. Moving the workspace label into the
resolved record instead of the URI would require a record
fetch on every boundary check, which either introduces a
per-call round trip on the enforcement path or requires an
aggressive local cache that reintroduces workspace-membership
consistency questions. The prefix-detection design is a
real operational property of the AGRP composition, not a
free choice, and it needs the workspace to be visible in
the URI to work.

The concession is explicit: everything about workspaces
that behaves as mutable governance metadata (membership
churn as agents move between workspaces, dual-workspace
membership per \S\ref{sec:proposals:owa:open} Q3,
workspace-level access lists, workspace deprecation
timestamps) lives in the resolved endpoint record and the
workspace catalogue, not in the URI. What the URI carries
is the workspace's role as a policy sub-scope of the
organizational trust atom, which is stable identity-defining
information (in the sense of the trust boundary the agent
operates within), even though workspace membership itself
is mutable.

We name this trade-off because a future revision of OWA,
or a related scheme, may reasonably choose differently:
if AGRP's prefix-detection property is not required
(e.g., if boundary detection is delegated to a mediator
that already fetches records), the trust-atom argument
alone does not compel keeping workspace in the URI, and
moving it to the resolved record would leave a cleaner
identity/metadata separation.

OWA is designed to coexist with the existing identifier
schemes rather than to replace them.

\paragraph{A2A AgentCard.}
The A2A AgentCard schema gains an optional
\texttt{agent\_uri} field carrying the agent's OWA URI.
Existing AgentCard URIs at \texttt{/.well-known/agent.json}
continue to work as transport endpoints; the OWA URI
sits alongside them as the stable identity. Rodriguez
proposes the same field for
\texttt{agent://}~\cite{agent_uri_scheme}; OWA reuses the
proposal verbatim.

\paragraph{MCP.}
MCP servers identified by stdio or HTTP transport endpoints
gain an optional OWA URI in the server manifest. The
workspace tier maps naturally onto MCP's emerging
tenant-scoping work in the AAIF working group
(\S\ref{sec:governance:foundations}).

\paragraph{ANP.}
ANP's DID-based identity resolves to a DID document;
the DID document gains an optional \texttt{owa\_uri} field
binding the DID to the OWA URI. The two identifiers
cross-reference: the OWA URI carries the workspace and
capability semantics; the DID carries the decentralized
verifiable identity. The bridge between them
(P8 of \S\ref{sec:openproblems:bridging}) is exactly what
the ANP-to-A2A composition that the rubric matrix most
clearly recommends needs.

\paragraph{ANS.}
ANS identifiers translate to OWA by mapping
\texttt{Provider} to \texttt{<org>} and inferring
workspace from a provider-side convention until the
provider migrates to native OWA. The migration is
incremental: ANS providers that adopt OWA can publish a
bridge entry that maps the legacy ANS name to the OWA
URI.

\subsubsection{Operational integration: IAM alias and email triggers}
\label{sec:proposals:owa:operational}

Two operational integration points warrant first-class
treatment because they are the cases where the
email-shape pattern (\S\ref{sec:proposals:owa:related})
genuinely outperforms the path form. OWA accommodates both
through optional, normatively-specified aliases that do
\emph{not} replace the canonical URI identifier.

\paragraph{IAM-operator-friendly alias.}
Compliance, operations, and security teams already know
how to govern email-shaped identities. Service
accounts, distribution lists, and automated workflows are
routine units of IAM management. OWA optionally defines a
canonical email-shaped alias of the form
\texttt{<agent-id>.<capability-slug>@<workspace>.<org>},
derived mechanically from the URI's components. For
example, the URI

\begin{quote}
\texttt{agent+owa://acme.com/w/finance/c/workflow/approval/}\\
\texttt{invoice/a/agent\_01h455vb4pex5vsknk084sn02q}
\end{quote}

\noindent has the canonical alias

\begin{quote}
\texttt{agent\_01h455vb4pex5vsknk084sn02q.}\\
\texttt{workflow-approval-invoice@finance.acme.com}
\end{quote}

The alias is one-way (parsing the alias mechanically
recovers the URI components, with the capability slug
hyphenated for local-part compatibility) and does
\emph{not} imply a deliverable mailbox. Some tooling
expects email-shaped identifiers: SIEM dashboards,
IAM consoles, audit reports, and ticket-routing systems.
Such tooling can render the alias without operators
needing to learn URI syntax. The form is intentionally consistent with
Google Cloud Service Accounts and Microsoft Entra UPNs
at the operator level~\cite{gcp_service_accounts,
entra_service_principals}, so an organisation migrating
from one of those IAM systems can preserve its operator
workflows even after adopting OWA URIs as the canonical
agent identifier. The alias is a presentational
affordance, not a parallel identity.

\paragraph{Trigger-by-real-email.}
A non-trivial fraction of production agent invocations
begin with a real email: a customer emails a support
inbox, a system sends a daily report, a CI hook emails an
alert, a calendar service sends a meeting reminder. OWA
optionally defines a trigger-email mapping that captures
this pattern without coupling the inter-agent protocol to
SMTP transport. An organisation publishes at
\texttt{https://<org>/.well-known/agent-triggers.json}
a map from email-address patterns to OWA URIs:

\begin{verbatim}
{
  "triggers": [
    { "pattern": "support@acme.com",
      "agent": "agent+owa://acme.com/w/support/c/.../a/triage-01" },
    { "pattern": "billing-report-*@finance.acme.com",
      "agent": "agent+owa://acme.com/w/finance/c/.../a/billing-01" }
  ]
}
\end{verbatim}

The organisation's mail infrastructure routes matching
inbound mail to the corresponding agent's A2A ingress
endpoint, packaging the email as an A2A
\texttt{Task} with the agent's OWA URI as the receiver.
The trigger email is a deliverable mailbox; the OWA URI
is the canonical agent identity; the two are bound by
org-controlled routing rather than by the protocol
specification. This composition separates the
trigger-channel concern from the inter-agent protocol
proper, gives the human-to-agent boundary a familiar
on-ramp, and preserves the path form's expressiveness
for cross-organizational inter-agent traffic. AGRP
applies unchanged: messages produced by the
mail-trigger path enter the agent's A2A surface and are
subject to the same egress policy as any other inbound
message.

\paragraph{What this gets the proposal.}
The two aliases together close the legitimate operational
gap that the email-shape pattern surfaces: organizations
whose existing IAM, audit, and trigger infrastructure is
email-shaped can adopt OWA without abandoning that
infrastructure. The path form remains the canonical
identifier and carries the capability semantics
(\S\ref{sec:proposals:owa:syntax}); the email-shaped
alias is a deterministic projection for tooling that
prefers email-syntax; the trigger-email map is a
routing-layer bridge that has no normative weight beyond
the org publishing it. OWA does not require either
alias to be configured. The point is to make adoption
possible without forcing a green-field rebuild of
identity infrastructure.

\subsubsection{Migration story}
\label{sec:proposals:owa:migration}

Three concerns determine whether OWA can be adopted
without re-architecting existing inter-agent deployments.

\paragraph{Organizations that have not normalized workspace names.}
The most common deployment state at the 2026-Q1 cutoff
is that organizations have no formal workspace
naming convention. OWA accommodates this with a
\texttt{default} workspace that an organization can use
until its naming convention stabilises:
\texttt{agent+owa://acme.com/w/default/c/.../a/...}. The
\texttt{default} workspace participates in AGRP
boundary detection as a \texttt{same\_workspace} destination; an
organization that later carves out workspaces from
\texttt{default} can do so by extending its workspace
catalogue without breaking existing references.

\paragraph{Workspace renaming.}
Workspaces are organizational concepts and can be
renamed (e.g., \texttt{customer-support} renamed to
\texttt{revenue-ops}). The OWA URI is stable against the
\texttt{<agent-id>} but not against the
\texttt{<workspace>} segment; an organization that
renames a workspace SHOULD publish an alias record at
the workspace catalogue mapping the old name to the new
name. Aliased lookups resolve transparently for a
migration window the organization specifies.

\paragraph{Coexistence with existing AgentCard URIs.}
Existing A2A AgentCard URIs continue to work as transport
endpoints during the migration. The OWA URI is a
\emph{name} that resolves to a transport endpoint; the
existing AgentCard URI is one such endpoint. Adopters
can deploy OWA without breaking any A2A client; clients
that don't yet speak OWA continue to use the AgentCard
URI directly.

\paragraph{Coexistence with existing \texttt{agent://} URIs.}
Organizations that have already deployed
\texttt{agent://} URIs~\cite{agent_uri_scheme} continue to
resolve them under the bare \texttt{agent://} scheme.
OWA URIs use the distinct \texttt{agent+owa://} scheme
suffix, so a resolver receiving an \texttt{agent://} URI
knows unambiguously to apply the three-tier grammar, and a
resolver receiving an \texttt{agent+owa://} URI applies
the four-tier OWA grammar. A non-OWA-aware
\texttt{agent://} parser presented with an
\texttt{agent+owa://} URI rejects the URI on the scheme
comparison rather than silently mis-parsing the path into
a wrong canonical name (the failure mode
\S\ref{sec:proposals:owa:related} identifies). Adoption is
opt-in per identifier: an organization can issue new
identifiers under \texttt{agent+owa://} while continuing
to honor previously issued \texttt{agent://} identifiers
for their agents' lifetimes.

\subsubsection{Open questions and caveats}
\label{sec:proposals:owa:open}

\paragraph{Q1: Workspace catalogue governance.}
The workspace catalogue at
\texttt{/.well-known/workspaces.json} is an
organization-controlled resource. The schema for the
catalogue (workspace list, parent-workspace relations,
deprecation timestamps) is not specified in OWA itself.
A future revision should specify the catalogue schema
normatively; for the 2026-Q1 cutoff, the field carries
the workspace name as a free-form string within the
URI regex.

\paragraph{Q2: Nested workspaces.}
Large organizations have hierarchical workspaces
(e.g., a \texttt{finance} workspace with sub-workspaces
\texttt{accounts-payable} and \texttt{accounts-receivable}).
OWA's \texttt{<workspace>} segment is currently flat. A
future revision should consider whether nested
workspaces should be encoded as
\texttt{/w/finance/accounts-payable/} or whether they
should be flattened with a dotted convention. The first
is more naturally extensible; the second is closer to
existing OAuth scope conventions.

\paragraph{Q3: Cross-organizational workspace identity.}
Two organizations may both have a \texttt{finance}
workspace, and a federated deployment may need to
distinguish between them. OWA disambiguates these via
the \texttt{<org>} segment; the open question is whether
a single agent can be a member of two workspaces in two
organizations simultaneously, and if so, how its URI
expresses the dual membership.

\paragraph{Q4: Compact-record pointer convention.}
Composing OWA with a lean-index registry such as
NANDA's AgentAddr~\cite{nanda_index}
(\S\ref{sec:proposals:owa:nanda}) leaves the anchor
record's URL fields reachable only by indirection: the
stored name is a one-way SHA-256 digest from which no
URL can be reconstructed. The pointer convention (a
digest-keyed registry bucket, a trust-root-derived
AgentFacts location, or another indirection mechanism)
is unspecified on both sides and needs co-design with
the registry's replication and caching model.

\paragraph{Q5: What the name denotes.}
The hardest open question is not syntax. It is what the
identifier refers to. Suppose the same agent code runs at
two URLs. The name may denote the software artifact, one
logical service with two deployments, or two agents.
ADR-0015 answers part of this for \texttt{urn:air}: its
rationale requires that relocating a workload not break
routing or caches, so the identifier denotes a logical
artifact rather than a deployment~\cite{ai_catalog}. OWA
inherits that reading, because the workspace tier is an
organizational scope and not a host. The unanswered half is
caller-dependent behaviour. Microsoft raised it when we
contacted it about this survey, and we did not find it
ourselves. If an endpoint exposes
different skills to different callers, it is undecided
whether that is one agent observed in two contexts or two
agents sharing a host. Nothing in the identifier grammar
decides it, and the AI Catalog work declares
identity-varying card generation out of scope and leaves it
to individual protocols. OWA does not resolve this either,
and we flag it rather than assume the logical reading
settles it. A candidate answer is that
identity is fixed by the publisher's declaration and
caller-dependent views are authorization outcomes rather
than distinct agents, which would tie the naming layer to
Axis 3 and make the denotation rule a security property.
Establishing that is future work.

\subsubsection{Composition with lean-index registries: the NANDA AgentAddr case}
\label{sec:proposals:owa:nanda}

The resolution machinery of
\S\ref{sec:proposals:owa:resolution} assumes some
registry substrate without mandating one. The closest
published design for that substrate at deployment scale
is the NANDA Index~\cite{nanda_index}: a distributed
registry for an ``Internet of AI Agents'' whose anchor
tier holds Ed25519-signed \textbf{AgentAddr} records
constrained to at most 120 bytes. Rich metadata
(capabilities, endpoints, telemetry) is pushed one tier
down into \textbf{AgentFacts}, W3C Verifiable Credential
objects. Federation is handled by a ``registry quilt'' of
CRDT-replicated, cross-signed registries. The design is
the natural complement to this section's proposal: NANDA
specifies the registry infrastructure that
\texttt{agent://} and OWA leave open, and
\texttt{agent://} specifies the naming layer that NANDA
leaves ad hoc. This subsection makes the composition
concrete, including where it does not compose.

\paragraph{The naming layer is NANDA's underspecified point.}
An AgentAddr record carries an \texttt{agent\_id}
(a \texttt{nanda:}-prefixed generic UUID) and an
\texttt{agent\_name} (a URN-style alias). The published
examples span several syntaxes
(\texttt{agent:Company:TranslationAssistant},
\texttt{urn:agent:salesforce:starbucks}) with no grammar,
no normalization rule, no capability semantics, and no
trust scoping. The consequence is structural: because
capabilities exist only in the AgentFacts tier, the index
tier cannot answer a capability-scoped query (``agents
under acme.com that can approve invoices'') without
fetching and parsing the heavy layer for each candidate.
Two records naming the same agent under different alias
syntaxes cannot be identified as equal. These are
precisely the properties the \texttt{agent://} naming
discipline provides~\cite{agent_uri_scheme}: a normative
ABNF, a canonical form, prefix-matchable capability
paths, and a DNS-anchored trust root. The two designs are
complementary layers, not competitors.

\paragraph{Aligning upstream does not close the naming gap.}
The reference implementation has since begun aligning its
index record to an external schema. Migration~010 adds
\texttt{version} and \texttt{trust\_manifest} columns for
parity with the AI Catalog \texttt{CatalogEntry}
schema~\cite{nanda_index_impl}, which makes AI
Catalog~\cite{ai_catalog} an upstream standard that NANDA
consumes rather than the reverse
(\S\ref{sec:proposals:owa:related}). The alignment sits at
the trust-metadata layer rather than the identifier layer:
the two columns carry a SemVer string and a signable
manifest of identity, attestations, provenance, and a
detached signature, and both are added to the organization
row that already scopes the record. Nothing in the
migration constrains \texttt{agent\_name}. Should the
identifier layer follow, AI Catalog's
\texttt{urn:air:\{publisher\}:\{namespace\}:\{name\}}
supplies the DNS-anchored publisher segment the published
examples lacked, but it remains an open text format with
no normative ABNF, no canonical form, and no stated
prefix-matchability, leaving name equality and
capability-scoped indexing at the anchor tier unavailable.
Being variable-length text of unbounded size, it also
leaves the arithmetic of the following paragraph
unchanged.

\paragraph{The 120-byte budget forces derivation, not storage.}
The constraint that makes the composition interesting is
arithmetic, and it follows from the record layout as
described in~\cite{nanda_index} rather than from any
deployed encoding. In the raw binary form that design
assumes, the fixed fields of an AgentAddr consume 84
bytes: the Ed25519 signature is 64, the UUID payload is
16, and the TTL is 4. That leaves
roughly 36 bytes for the \texttt{agent\_name} and three
URL fields (\texttt{primary\_facts\_url},
\texttt{private\_facts\_url},
\texttt{adaptive\_resolver\_url}). A single typical HTTPS
URL runs 60 bytes or more, so no literal encoding of even
one field fits. The budget therefore rules out storing
\emph{any} name or pointer literally and forces every
variable-length field through some derivation or
indirection. Once that is accepted, the length of the
canonical URI string stops mattering to the anchor tier
entirely, because what is stored is a fixed-width digest.

\paragraph{The budget is specified, not yet enforced.}
The arithmetic above is a property of the published
design rather than of the current reference
implementation, whose index record is a relational row
with unbounded \texttt{JSONB} columns, scoped one record
per publishing organization rather than per agent, and
carrying no anchor-tier byte
constraint~\cite{nanda_index_impl}. The figures should
therefore be read as what a compliant anchor tier must
satisfy, not as a measurement of deployed code. The
composition is unaffected either way: a derived,
fixed-width name is what makes the record fit when the
budget is enforced, and costs nothing when it is not.
The implementation is also moving toward the published
design rather than away from it, having since gated
activation on a DNS TXT challenge proving control of the
claimed domain, on the stated reasoning that the domain
is the namespace the index is keyed on, and added a
signed trust manifest carrying identity, attestations,
and provenance.

\paragraph{Field-by-field mapping.}
Table~\ref{tab:nanda-field-mapping} maps the five AgentAddr
fields onto the naming components, one row per field.

\begin{table}[h]
  \centering
  \footnotesize
  \begin{tabular}{p{2.4cm}p{3.4cm}p{4.2cm}}
    \toprule
    AgentAddr field & NANDA as published & agent:// / OWA contribution \\
    \midrule
    \texttt{agent\_id}
      & \texttt{nanda:} + generic UUID (random, string-encoded)
      & TypeID: the 16-byte binary UUIDv7 payload, time-sortable for range scans and shard routing \\
    \texttt{agent\_name}
      & ad-hoc URN alias, no grammar
      & the 32-byte SHA-256 digest of the canonical URI's identity-defining components (the key of \S\ref{sec:proposals:owa:resolution}); verifiable against the full URI in AgentFacts \\
    URL fields
      & literal URLs, individually over budget
      & not derivable from the URI; require a pointer convention (Q4, \S\ref{sec:proposals:owa:open}) \\
    \texttt{signature}
      & Ed25519 over the record
      & composes unchanged; capability claims additionally carry PASETO v4.public attestation bound by the coverage rule of~\cite{agent_uri_scheme} \\
    \texttt{ttl}
      & seconds
      & unchanged \\
    \bottomrule
  \end{tabular}
  \caption{Field-by-field mapping between NANDA AgentAddr records and agent:// / OWA naming components.}
  \label{tab:nanda-field-mapping}
\end{table}

Two of the five rows are strict improvements with no
interface change. A random UUID gives a sharded global
index nothing to sort or range-scan on; a UUIDv7 payload
is time-ordered at identical width. Replacing an
ungoverned alias string with the canonical digest gives
the anchor tier name equality, spoof resistance
(recomputing the digest from a claimed URI either matches
the stored name or does not), and capability-scoped
indexing at the tier where discovery runs.

\paragraph{What derives and what does not.}
The digest is an index key, not a compression scheme:
SHA-256 is one-way, so no URL field can be reconstructed
from the stored name. The three URL fields are genuinely
dynamic per deployment and must be reached through
indirection (for instance, a registry bucket keyed by the
stored digest whose record carries the pointers, or an
AgentFacts document fetched by a convention derived from
the trust root's DNS name). Either convention keeps the
anchor record within budget (84 fixed + 32 digest = 116
bytes) at the cost of one indirection on the cold path,
the same cost the AgentFacts tier already imposes. The
pointer convention itself is not specified by
\texttt{agent://}, OWA, or the NANDA publications
reviewed; it is the genuine co-design problem between the
naming layer and the index, recorded as Q4 in
\S\ref{sec:proposals:owa:open}.

\paragraph{The workspace tier composes without touching the anchor record.}
The asymmetry of \S\ref{sec:proposals:owa:resolution},
in which the workspace label is resolution-time policy
metadata rather than an input to key derivation, pays off
here: an \texttt{agent+owa://} URI and its underlying
identity hash to the same anchor-tier name. A NANDA
deployment adopting OWA changes nothing at 120-byte
scale. The workspace label, its attestation, and its
catalogue entry live in the AgentFacts tier alongside
other rich metadata, where NANDA already places mutable
governance facts. This is the same identity/metadata
split defended in \S\ref{sec:proposals:owa:tradeoff},
arrived at independently by both designs.

\paragraph{Scope of this composition.}
\texttt{agent://} and OWA occupy the
\texttt{agent\_name} and \texttt{agent\_id} fields and
supply deterministic index keys; they do not replace
the registry quilt, the AgentFacts credential schema,
the TTL-driven resolution tiers, or NANDA's
privacy-path machinery. Nothing above should be read as a
claim that a URI scheme substitutes for registry
federation. The privacy-preserving discovery guarantees
that \texttt{agent://} lists as future work remain
NANDA's contribution, not ours. A broader survey of the
registry-solution landscape, including DNS-based,
DID-based, and distributed-index approaches, is provided
by Singh et al.~\cite{agent_registry_evolution}.

\subsection{AGRP Guardrails Envelope (Contribution C6)}
\label{sec:proposals:agrp}

\subsubsection{Problem statement}
\label{sec:proposals:agrp:problem}

Three of the six open security problems of
\S\ref{sec:security:open} concern inter-agent message
content rather than identity or transport:
O3 (prompt-injection mitigation as spec-level requirement),
O5 (contextual integrity for conversation-history
forwarding), and O6 (side-channel metadata privacy). The
healthcare and finance disclosure gap of
\S\ref{sec:adoption:healthcare} to \ref{sec:adoption:finance}
is the institutional analogue: organizations in those
verticals deploy intra-organizational inter-agent
coordination but resist cross-organizational deployment
because the contemporary protocols do not give them
spec-level control over what content crosses an
organizational boundary.

The deployment-layer response to this gap is the egress
proxy. The most-developed example is
Pipelock~\cite{pipelock}, which provides an out-of-band
proxy with content inspection (DLP, prompt-injection, and
MCP tool policies), an RFC-8941 Mediation Envelope, a
hash-chained Flight Recorder audit log, and
mediator-signed Action Receipts with a standalone offline
verifier CLI whose verification key is obtained
out-of-band and pinned before receipt validation. Pipelock is the
state of the art for the egress-proxy pattern; the survey
pins its citation to the \texttt{v3.0.0} tag because
Pipelock's pattern-catalogue counts grow release-to-release
and any concrete number cited here would date faster
than the survey's own review cycle. The survey's
companion rubric artefact (\S\ref{sec:conclusion:artifact})
is the mechanism by which successor tags stay tracked
against the same rubric.

Pipelock's own documentation flags the principal
limitation of the deployment-layer approach: ``With the
sandbox enabled or the host/cluster containment topology
enforced, [the proxy] adds an OS or deployment boundary
on top of content inspection. Direct egress still has to
be blocked by that boundary for non-cooperative tools
that ignore proxy settings''~\cite{pipelock}. Stated more
generally: a deployment-layer mediator is
\emph{discoverable} and \emph{bypassable} for an
attacker who can manipulate environment variables, the
network stack, or the agent's own configuration. The
defence is to push the enforcement up to the message
layer, where the bypass is no longer available because
the receiver checks the envelope.

AGRP is that message-layer envelope: a normative
extension to A2A that wraps every cross-boundary message
with classification labels, a tamper-evident redaction
log, and attestation that an egress filter ran. The
attestation is bound to the sender's OWA identity
(Contribution C5) rather than to a mediator's identity,
making the egress-filtering property a property of the
message itself rather than of the deployment that
produced it.

\subsubsection{Related work and differentiation}
\label{sec:proposals:agrp:related}

Table~\ref{tab:agrp-differentiation} compares AGRP
against ten prior works in the guardrails space.

\begin{table*}[t]
  \centering
  \caption{AGRP differentiation matrix. The seven-property
    argument: no prior system simultaneously provides all
    seven of continuous redaction, filter attestation,
    receiver verifiability, naming-derived boundary
    detection, multi-backend attestation composition,
    cross-vendor verifiability, and normative spec-level
    requirement.}
  \label{tab:agrp-differentiation}
  \footnotesize
  \resizebox{\textwidth}{!}{%
  \begin{tabular}{p{2.7cm}p{2.6cm}cccccc p{2.5cm}}
    \toprule
    Prior work & Mechanism & Continuous? & Filter attestation? & Receiver verifies? & Boundary from naming? & Cross-vendor? & Normative? & AGRP improvement \\
    \midrule
    Improving A2A~\cite{a2a_weaknesses}
      & \texttt{USER\_CONSENT} + tunnels
      & transactional & no & no & no & A2A & partial
      & continuous; addresses silent leakage \\
    Pipelock~\cite{pipelock}
      & egress proxy + content-inspection catalogue + RFC-8941 + signed receipts
      & yes (proxy) & \textbf{yes (mediator)} & \textbf{yes (offline CLI)} & no & per-deployment & no
      & AGRP carries Pipelock receipts under a normative Tier~2 semantics; binds to sender identity, not mediator \\
    Ruflo PII pipeline~\cite{ruflo}
      & 14-type detect + BLOCK/REDACT/HASH/PASS
      & yes (per node) & no & trust-score-based & no & one orchestrator & no
      & filter-attested vs. trust-scored (preventive vs. reactive) \\
    NeMo Guardrails
      & programmable rails
      & per agent & no & no & no & no & no
      & inter-agent; receiver verifies filter ran at sender \\
    Galileo / FutureAGI / Maxim
      & dual-stage I/O validation
      & per orchestrator & no & no & no & no & no
      & cross-orchestrator vendor boundary \\
    OpenHarness~\cite{openharness}
      & Pre/PostToolUse hooks + permission modes
      & per tool call & no & no & no & per-deployment & no
      & protocol-layer; OpenHarness is harness-layer \\
    Semgrep~\cite{semgrep_a2a}
      & capability-based authz recommendation
      & n/a & no & no & no & n/a & no
      & normative envelope, not a recommendation \\
    SAGA~\cite{saga2026}
      & TEE-anchored gateway
      & yes (gateway) & \textbf{yes (TEE quote)} & yes & no & per-deployment & no
      & AGRP $\supset$ SAGA's TEE quote as one backend \\
    Proof-of-guardrail~\cite{proof_of_guardrail}
      & agent + named guardrail in a TEE, per-response attestation
      & yes (per response) & \textbf{yes (TEE)} & \textbf{yes (offline)} & no & one TEE vendor & no
      & inter-agent rather than user-to-agent; AGRP carries the quote as its Tier~2 hardware path and normalizes it across vendors \\
    A2A v0.3+ AgentCard signing
      & optional card signing
      & per agent & yes (identity) & yes (identity) & no & per-A2A & optional
      & attests content filter, not just identity \\
    \bottomrule
  \end{tabular}%
  }
\end{table*}

Two cells of the table carry the differentiation argument
most heavily. We give each a dedicated paragraph. The
proof-of-guardrail row~\cite{proof_of_guardrail} carries a
third argument that is about assurance rather than
positioning, so \S\ref{sec:proposals:agrp:tee} takes it up
after the tiers are defined.

\paragraph{The Pipelock relationship: composition, not competition.}
Pipelock is the closest prior work and the most
operationally mature implementation in this space. The
naive reading of AGRP is that it is ``Pipelock at the
protocol layer'' and that the right deployment posture is
to use Pipelock without AGRP. Three architectural
differences argue against this reading.

First, \emph{where the filter lives}. Pipelock runs in a
proxy process external to the agent. AGRP allows the
filter to run anywhere: in-process, in a sidecar, in a
Pipelock mediator, or in a SAGA TEE gateway. The only
condition is that the attestation in the AGRP envelope
binds to the sender's OWA URI. This decouples the policy from the
deployment.

Second, \emph{who attests}. Pipelock receipts are signed
by the mediator (the Pipelock process, a boundary the
sender does not control); AGRP envelopes are signed by
the sender's identity, with the mediator's receipt
nested inside the envelope's \texttt{mediator\_receipt}
field. A receiver verifying an AGRP envelope
independently validates both layers: the
sender-identity signature confirms the message came
from the named agent; the nested mediator receipt
confirms an independent filter actually ran. The two
attestation chains are orthogonal, and AGRP's tier
framing (\S\ref{sec:proposals:agrp:tiers}) makes the
distinction explicit: a Pipelock receipt inside the
envelope elevates the assurance level from Tier~1
(sender-signed claim) to Tier~2 (independently-signed
evidence), and the tightenings of that subsection are
what keep the elevation honest.

Third, \emph{cross-vendor composition}. Pipelock requires
the proxy to be co-located with the agent (via
\texttt{HTTPS\_PROXY} or sidecar). AGRP is a message-level
invariant any A2A speaker can carry, so an agent at
organization X can verify the receipt from organization
Y without organization Y running Pipelock. Organization Y
only needs to run \emph{some} filter that produces an
AGRP-compatible attestation. Pipelock-as-mediator under
AGRP is one such filter; an in-process filter on
organization Y's agent is another; a SAGA TEE gateway is
a third. Cross-vendor composition is the property
Pipelock cannot provide on its own.

The composition pattern (\S\ref{sec:proposals:agrp:composition})
makes the relationship explicit: AGRP envelope
$\supseteq$ Pipelock receipt. A
Pipelock-deployment plus AGRP-protocol pair gives both
the proxy's network-level enforcement \emph{and} the
protocol-level cross-vendor invariant; the deployment is
defence in depth rather than one substituting for the
other.

\paragraph{The Ruflo relationship: preventive vs. reactive.}
Ruflo's 14-type PII detection pipeline applied per
trust-level (BLOCK / REDACT / HASH / PASS) is the
operational equivalent of AGRP's classification labels
and redaction actions~\cite{ruflo}. The architectural
difference is the attestation model. Ruflo computes a
trust score from behavioural signals (success rate,
uptime, threat detection, integrity) and downgrades a
misbehaving agent's trust level. The downgrade is
\emph{reactive}: by definition, a downgraded agent
has already leaked once. AGRP's filter attestation is
\emph{preventive}: the filter must run \emph{before} the
message is signed, and the receiver fails closed if the
attestation does not validate. Trust-scoring and
filter-attestation are not mutually exclusive. A
deployment can use Ruflo's trust score to choose which
filter rule-set applies to a message, then AGRP-attest
that the chosen rule-set ran. They compose in the same
way Pipelock does.

\subsubsection{The AGRP envelope: schema and semantics}
\label{sec:proposals:agrp:schema}

The AGRP envelope wraps every cross-boundary A2A message
with a normative structured-field block adjacent to the
A2A message body. The envelope is JSON-encoded at the
A2A transport boundary; the sideband-metadata equivalent
on HTTP transports uses RFC 8941 Structured Field
Values~\cite{rfc8941}, matching Pipelock's Mediation
Envelope choice for compatibility.

\begin{verbatim}
{
  "agrp_version": "1.0",
  "sender":    "agent+owa://acme.com/w/finance/c/.../a/billing-01",
  "recipient": "agent+owa://acme.com/w/support/c/.../a/chat-02",

  "boundary": {
    "kind": "cross_workspace",
    "egress_policy_id":  "acme/finance.outbound.v3",
    "ingress_policy_id": "acme/support.accept.v2"
  },

  "classification": {
    "tier": "internal-confidential",
    "data_classes": ["financial_record"],
    "applied_labels": ["pii:redacted", "secret:blocked",
                       "env_var:blocked"]
  },

  "redaction_log": [
    { "jsonpath": "$.message.text",
      "rule": "redact_ssn",
      "count": 1,
      "fingerprint": "sha256:5f2..." },
    { "jsonpath": "$.message.text",
      "rule": "block_stripe_key",
      "count": 1,
      "action": "blocked" },
    { "jsonpath": "$.parts[2].data",
      "rule": "block_env_var",
      "count": 3,
      "action": "stripped" }
  ],

  "attestation": {
    "filter_id":          "acme-guardrails-v3.2",
    "filter_measurement": "sha256:1d4f...64-hex-chars...",
    "signing_key_id":     "owa-key-2026q3",
    "assurance_tier":     1,
    "evidence_status":    "self_attested_claim",
    "filter_signature":   "ed25519:base64url...",
    "mediator_receipt":   "...",
    "tee_quote":          "..."
  },

  "payload": { /* A2A message with redactions applied */ }
}
\end{verbatim}

The fields carry the following semantics. The
\texttt{sender} and \texttt{recipient} fields are OWA
URIs (Contribution C5). The \texttt{boundary} field
records the boundary kind derived from prefix comparison
of the sender and recipient URIs
(\S\ref{sec:proposals:agrp:boundary}) and the policy
identifiers that applied at egress and that the
recipient should apply at ingress. The
\texttt{classification} field records the
sensitivity tier and the labels that the egress filter
identified or applied; the label vocabulary is open
(an organization can extend it) but the protocol
specifies a minimum set of nine well-known labels
(\texttt{pii}, \texttt{secret}, \texttt{env\_var},
\texttt{financial}, \texttt{phi}, \texttt{credentials},
\texttt{token}, \texttt{contract}, \texttt{audit}).

The \texttt{redaction\_log} is the most operationally
consequential field. Each entry records the JSONPath at
which a rule fired, the rule identifier, the number of
matches, and either a SHA-256 fingerprint of the redacted
content (allowing later correlation when an authorised
verifier has the original or an independently trusted
receipt binds the same commitment) or an action marker
(\texttt{blocked}, \texttt{stripped}) for content that
was removed entirely. The sender signature makes the
declared fingerprint log tamper-evident in transit. A
fingerprint alone neither proves that the committed
plaintext matched the claimed rule nor proves that the
sender reported every match. At Tier~1 both are sender
claims. Evidence that a filter actually produced a
complete log requires a Tier~2 receipt whose independently
verified statement binds the payload, filter measurement,
and log commitment.

The \texttt{attestation} block records a signed
enforcement claim and, when present, independently
verifiable evidence. The \texttt{filter\_id} and
\texttt{filter\_measurement} identify the filter binary;
the \texttt{filter\_signature} is signed by the sender's
OWA-bound key and authenticates the sender's claim but does
not by itself prove execution; the \texttt{mediator\_receipt} field is
optional and carries an independent-mediator receipt
(Pipelock-style Action Receipt~\cite{pipelock}, or an
equivalent) when the filter ran outside the sender's
trust boundary. The field name is AGRP's, and the
receipt itself follows the mediator's own on-wire
format. Pipelock v3.0.0, for instance, structures its
own Mediation Envelope as an RFC-8941 dictionary
carried in a \texttt{Pipelock-Mediation} HTTP header
rather than under this field name. The
\texttt{tee\_quote} field is optional and carries a
TEE-anchored quote (SAGA-style~\cite{saga2026}) when the
filter ran inside vendor-attested hardware. The
combination of which of these three signatures are
present determines the envelope's attestation tier
(\S\ref{sec:proposals:agrp:tiers}); the schema also
reserves an optional \texttt{log\_anchor} sub-field
inside the attestation block for the append-only-log
inclusion proof described in the same subsection, kept
outside the required schema at the cutoff so that
early implementations are not forced to run a
transparency log to be conformant.

Both AGRP and ACSP use the same signing profile in the
reference code: RFC~8785 JSON Canonicalization is applied
to the complete envelope with only the protocol signature
field removed, and Ed25519 signs those bytes. The signature
is base64url encoded with the \texttt{ed25519:} prefix. The
\texttt{signing\_key\_id} selects an OWA-bound public key
provisioned independently of the envelope. The signed
\texttt{assurance\_tier} and \texttt{evidence\_status}
fields prevent an optional opaque receipt from silently
upgrading a Tier~1 claim.

\subsubsection{Boundary detection from OWA URIs}
\label{sec:proposals:agrp:boundary}

The boundary kind in the envelope's \texttt{boundary}
field is derived by prefix comparison of the sender and
recipient OWA URIs (Contribution C5). The canonical
on-wire enum uses snake case:
\texttt{same\_workspace}, \texttt{cross\_workspace}, and
\texttt{cross\_org}. The procedure is:

\begin{enumerate}
  \item Parse both URIs into their org / workspace /
    capability / agent-id components.
  \item If the \texttt{<org>} components differ, the
    boundary is \texttt{cross\_org}.
  \item Otherwise, if the \texttt{<workspace>} components
    differ, the boundary is \texttt{cross\_workspace}.
  \item Otherwise, the boundary is \texttt{same\_workspace}.
\end{enumerate}

The egress filter consults the sender's per-workspace
policy catalogue to find the egress rule-set that
applies to the boundary kind. AGRP is mandatory for
\texttt{cross\_workspace} and \texttt{cross\_org}
boundaries and optional for \texttt{same\_workspace}
messages, on the operational principle that
within-workspace traffic is already inside a single
policy domain and the overhead of envelope construction
is not warranted unless an explicit per-message policy
demands it.

The coupling between C5 and C6 is intrinsic: AGRP
without OWA would need to specify its own boundary
notion (which would duplicate the workspace concept
under a different name), and OWA without AGRP would
provide identity scoping but no protocol-layer
enforcement of the policy boundaries the workspace
identifies. The same boundary detection serves ACSP
(\S\ref{sec:proposals:acsp:composition}); the
broader triplet argument is set out in
\S\ref{sec:proposals:codesign}.

\subsubsection{Attestation, fail-closed semantics, and receiver verification}
\label{sec:proposals:agrp:attestation}

The receiver's verification procedure on an inbound AGRP
envelope is:

\begin{enumerate}
  \item Verify the \texttt{filter\_signature} against the
    sender's OWA-bound public key. If invalid, fail
    closed.
  \item Check the \texttt{filter\_id} and
    \texttt{filter\_measurement} against the receiver's
    allow-list of approved filters for the
    boundary kind. If the filter is not on the
    allow-list, fail closed.
  \item Determine the envelope's attestation tier
    (\S\ref{sec:proposals:agrp:tiers}) from the
    presence and independence of the
    \texttt{mediator\_receipt} and \texttt{tee\_quote}
    fields. If the receiver's ingress policy requires
    Tier~2 for this boundary kind or classification
    tier and only a Tier~1 attestation is present, fail
    closed.
  \item Validate the \texttt{mediator\_receipt} if
    present, using the mediator's verifier key obtained
    out-of-band ahead of the message (a
    key delivered inside the envelope MUST NOT satisfy
    this step).
  \item Validate the \texttt{tee\_quote} if present
    against the TEE vendor's attestation
    infrastructure.
  \item If the envelope carries the optional
    log-anchored sub-tier
    (\S\ref{sec:proposals:agrp:tiers}) and the
    receiver's ingress policy requires it for this
    message, verify the inclusion proof against the
    log's published root; if the proof is missing or
    invalid, fail closed.
  \item Apply the receiver's ingress policy to the
    \texttt{classification} labels and the
    \texttt{boundary} field. If the labels exceed the
    ingress policy's accept list, fail closed.
  \item Validate the \texttt{redaction\_log} schema and
    commitment syntax. The approved
    \texttt{filter\_measurement}, not the list of rules
    that happened to fire, identifies the complete rule
    set. At Tier~1 an absent match remains a sender claim;
    at Tier~2 the independent receipt must bind the log
    commitment and filter measurement.
\end{enumerate}

The fail-closed discipline is normative: an AGRP-aware
receiver MUST refuse to process the payload if any of
the verification steps fail. A receiver that downgrades
to ``process anyway with a warning'' breaks the
protocol's guarantees and the receiver SHOULD be flagged
non-compliant by an AGRP conformance test suite.

The verification chain is independent of the sending
deployment. A receiver at organization Y can validate
the AGRP envelope from organization X without sharing
infrastructure. The verification depends only on the
sender's OWA-bound key and the receiver's allow-list of
approved filters. This is the property that Pipelock
deployment alone cannot give: Pipelock's receipts are
verifiable only by parties that share Pipelock's PKI;
AGRP envelopes are verifiable across organizational
boundaries through DNS and the OWA trust chain.

\subsubsection{Attestation assurance tiers}
\label{sec:proposals:agrp:tiers}

The envelope schema of
\S\ref{sec:proposals:agrp:schema} accommodates a
range of attestation backends, but a receiver making a
trust decision needs to know which backend produced the
attestation it just verified. AGRP therefore defines
two normative assurance tiers, drawing an explicit line
between an \emph{enforcement claim} and \emph{evidence of
enforcement}. The line is the same one drawn by the
Pipelock maintainers' analysis of self-attestation vs.\
independent
mediation~\cite{pipelab_independent_attestation}: a
sender signing ``a filter ran'' with its own key is a
claim by the party that would benefit from lying, and a
boundary the sender does not control signing what it
actually did is evidence. Most of the field blurs the
two; the tiers below draw the line so that a receiver
can weight the two cases differently in its ingress
policy.

\paragraph{Tier 1: self-attested (low assurance).}
The filter runs inside the sender's own trust boundary.
That means in the agent's process, in a sidecar the
agent's operator controls, or in any component whose
failure mode is indistinguishable from the sender saying
``trust me.'' The attestation is signed only by the
sender's OWA-bound key
(\S\ref{sec:proposals:agrp:schema}, the
\texttt{filter\_signature} field). The envelope's
\texttt{mediator\_receipt} and \texttt{tee\_quote}
fields are absent. Tier~1 is a claim about enforcement,
not enforcement, and receiver ingress policy SHOULD
reflect that: Tier~1 is appropriate for
\texttt{same\_workspace} traffic and for low-sensitivity
\texttt{cross\_workspace} traffic within a single
organization, but SHOULD NOT be accepted for
\texttt{cross\_org} boundaries or for classification
tiers above \texttt{internal-confidential}.

\paragraph{Tier 2: independently attested (high assurance).}
The filter runs at a boundary the sender does not
control, and the attestation is signed by that
independent boundary. That boundary is either a
mediator process the sender cannot tamper with, or a
TEE anchored in vendor-attested hardware. The envelope carries either a
\texttt{mediator\_receipt} (mediator path) or a
\texttt{tee\_quote} (hardware path), and the receiver
verifies the signature against a key published by the
independent boundary ahead of time. Tier~2 provides
evidence of enforcement, not merely a claim, and is the
minimum assurance level receivers SHOULD require for
\texttt{cross\_org} traffic and for classification tiers
that carry regulated content (\texttt{financial},
\texttt{phi}, \texttt{contract}, \texttt{credentials}).
The canonical Tier~2 mediator implementation at the
2026-Q1 cutoff is Pipelock v3.0.0, whose Action
Receipts are signed by the proxy process and verifiable
offline by a standalone CLI against a key published
ahead of the traffic. An implementer can
reproduce that example end-to-end using the Pipelab receipt
playground~\cite{pipelab_playground} without contacting
a server.

\paragraph{Two tightenings that keep the Tier~1/Tier~2 line un-gameable.}
The Tier~2 definition is worth stating precisely
because two apparently-equivalent constructions
degrade it back to Tier~1:

\begin{enumerate}
  \item \emph{Independence covers enforcement, not
    only signature.} A sender that runs the filter
    in-process and then hands the resulting log to an
    outside party for a co-signature has produced an
    attestation signed by an independent boundary of a
    filter that ran inside the sender's trust boundary.
    That is Tier~1 in a Tier~2 envelope. Tier~2
    requires the independent boundary both to
    \emph{perform} the filtering and to sign the
    receipt. An AGRP conformance test suite verifying
    that a receipt is Tier~2-eligible SHOULD check both
    conditions. The \texttt{filter\_measurement}
    field must identify a binary whose execution
    environment is under the same independent
    authority as the signing key.
  \item \emph{Verifier key must be published
    out-of-band, ahead of the receipt.} A
    ``verifiable-offline'' receipt is only verifiable
    if the receiver already has the verifier key by
    the time the receipt arrives. A key delivered
    inside the same bundle it is meant to verify is
    circular: the receiver has to trust the transport
    to trust the key to trust the receipt. Tier~2
    requires the mediator's (or TEE's) verifier key to
    have been published in a location the receiver
    consults independently of the message. Examples are
    the mediator's HTTPS-served key endpoint, a DNSSEC
    record under the mediator operator's zone, a
    transparency log, or (for the TEE path) the vendor
    attestation infrastructure. The tightening is a
    property of the \emph{receiver's} verification
    procedure, not a prohibition on the receipt bundle
    carrying any key material: a self-describing
    receipt may embed the signer's public key for
    identification without that key being the trust
    anchor. Pipelock v3.0.0 illustrates the shape the
    tightening should take. The signer's public key is
    not carried in the receipt body itself; the verifier
    and chain layer pin a trusted signer key obtained
    out-of-band and compare each incoming receipt's
    signature against the pinned key at validation time.
    Unpinned verification is labelled structural-only and
    exits non-zero unless the operator opts in. This
    preserves the out-of-band-key property by
    construction: the trust anchor lives at the verifier
    layer rather than inside the receipt.
\end{enumerate}

\paragraph{Sub-tier 2$^\star$: log-anchored (defence against mediator rewrite).}
A Tier~2 receipt gives the receiver evidence that the
independent boundary attested to a specific state of
filtering at a specific moment; it does not by itself
prevent the independent boundary from
\emph{retrospectively} re-emitting a different receipt
for the same moment after the fact. Some contexts do not
fully trust even the mediator: adversarial
audit, disputes between mediator operator and sender,
and regulator-driven forensics. For those contexts, AGRP
defines an optional
sub-tier in which the mediator commits each receipt to
an external append-only log (a certificate-transparency
style Merkle log, an attestation blockchain, a witness
network) whose inclusion proof is carried in the
envelope alongside the receipt. Receivers who require
the sub-tier verify the inclusion proof against the
log's published root, giving them a receipt the
mediator cannot rewrite. The append-only-log field is
optional in the AGRP envelope
(\S\ref{sec:proposals:agrp:schema}); a future revision
may promote it to normative for regulated verticals.

\subsubsection{Composition: mapping backends to assurance tiers}
\label{sec:proposals:agrp:composition}

AGRP composes with three classes of egress-filter
backend, and each maps cleanly to a tier defined in
\S\ref{sec:proposals:agrp:tiers}. The mapping, not the
backend inventory, is what a receiver's ingress policy
should be written against.

\paragraph{In-process filter (Tier~1).}
The simplest deployment runs the egress filter inside
the agent's own process: the runtime calls into a
filter library, the filter signs the envelope with the
agent's OWA-bound key, and the envelope is
transmitted. The threat model is honest about what
this provides: a compromised agent process can
forge a clean envelope. That is why the tier
framing (\S\ref{sec:proposals:agrp:tiers}) locates it
in Tier~1 rather than describing it as ``weak Tier~2.''
It is appropriate for \texttt{same\_workspace} traffic
and for low-sensitivity intra-org \texttt{cross\_workspace}
traffic; it is not appropriate for \texttt{cross\_org}
egress or regulated classifications.

\paragraph{Pipelock mediator (Tier~2, mediator path).}
A Pipelock-style proxy~\cite{pipelock} runs out-of-band
to the agent process. The agent's egress traffic flows
through the proxy; the proxy performs the content
inspection (the enforcement) and emits a signed Action
Receipt (the evidence). AGRP carries the receipt in its
\texttt{mediator\_receipt} field; the agent additionally
signs the outer envelope with its OWA-bound key. Note
that \texttt{mediator\_receipt} is an AGRP field name,
not the container Pipelock itself uses on the
wire.\footnote{Pipelock v3.0.0's own transport encodes
a Mediation Envelope as an RFC-8941 structured-field
dictionary carried in a \texttt{Pipelock-Mediation}
HTTP header (or a \texttt{com.pipelock/mediation} key
in MCP \texttt{\_meta}), and Pipelock's agent-to-agent
envelope references a receipt by content digest plus
\texttt{signer\_key} rather than nesting the receipt
bytes. AGRP's \texttt{mediator\_receipt} is the carrier
of the receipt on the AGRP path specifically; the
Pipelock transport shape is preserved on Pipelock's own
path.} The receiver verifies both layers independently.
The agent's signature confirms the named agent endorsed
the message. The mediator's receipt confirms an
independent filter actually ran. That receipt is signed
by a boundary the agent does not control. The receiver
verifies it against a trusted verifier key obtained
out-of-band, per the Pipelock verifier's pinning
discipline~\cite{pipelab_independent_attestation}. This is
the canonical Tier~2 mediator instantiation and the
reference point for the tier definition; readers who
want to reproduce the offline-verification property
end-to-end can pull a signed receipt and verify it
in-browser using the Pipelab
playground~\cite{pipelab_playground}.

\paragraph{SAGA TEE gateway (Tier~2, hardware path).}
A SAGA-style gateway~\cite{saga2026} runs the filter
inside a TEE (Intel SGX, AMD SEV-SNP, Apple Secure
Enclave). The TEE produces a vendor-attested quote that
the gateway emits as part of the envelope's
\texttt{tee\_quote} field, and the receiver verifies
the quote against the TEE vendor's attestation
infrastructure. That is again a verifier the receiver
consults independently of the sender. The hardware
path shares the Tier~2 assurance level with the
mediator path but shifts the trust root from the
mediator operator to the TEE vendor, which is the
right choice for host-operator-untrusted scenarios
(multi-tenant public infrastructure, adversarial hosting
environments) even though it is the most operationally
expensive of the three backends.

\paragraph{Combining backends.}
The backends compose. An organization that requires
the strictest posture might deploy a SAGA TEE gateway
as the outermost AGRP filter, with a Pipelock mediator
wrapped inside the TEE, with an in-process filter as
a fallback for \texttt{same\_workspace} traffic.
Backends may also be combined with the append-only-log
sub-tier (\S\ref{sec:proposals:agrp:tiers}) to give
receipts the mediator itself cannot rewrite. AGRP
specifies the envelope and the tier semantics; the
backend choice is a deployment decision constrained
only by the receiver's ingress-policy floor for the
message's boundary kind and classification tier.

\subsubsection{What AGRP guarantees without a TEE, and when
a TEE is the right answer}
\label{sec:proposals:agrp:tee}

The tier definitions above make the TEE one of two Tier~2
paths. A reader is entitled to ask the sharper question:
what does AGRP actually guarantee when no TEE is present,
and what is the case for reaching for one? The
proof-of-guardrail scheme~\cite{proof_of_guardrail} is the
right comparator, because it answers the adjacent question
with hardware and reports what that costs. It runs the agent
and a named open-source guardrail inside a TEE and emits a
TEE-signed attestation that the guardrail executed before
the response was released. A user verifies the attestation
offline, without contacting the developer. That is a
per-response execution proof, and it is strictly stronger
evidence than any signature the sender produces about
itself.

The comparison is best stated as three claims and who can
make them.

\begin{description}
  \item[Tier 1, no TEE, no mediator.] AGRP guarantees the
    \emph{integrity and attribution of a claim}, not the
    execution of a filter. A verifying receiver learns four
    things. It learns which OWA-bound identity signed the
    envelope. It learns which filter identifier and
    measurement that identity claims to have run. It learns
    which spans that identity says it redacted. And it learns
    that none of this was altered in transit.
    That is worth having, and it is not enforcement. A
    sender who never ran the filter can sign the same
    envelope. What the receiver gains is a
    \emph{non-repudiable} record. The false claim is signed,
    attributed, and durable. The failure therefore becomes an
    accountability question after the fact rather than an
    invisible one. Tier~1 is therefore appropriate exactly
    where post-hoc accountability is a sufficient deterrent,
    which is inside one organization's own trust boundary.
  \item[Tier 2, mediator path.] AGRP guarantees that a
    boundary the sender does not control observed the
    traffic and signed what it saw. The residual trust moves
    to the mediator operator. This is the right choice when
    the sender and receiver are different organizations but
    both accept a common mediator, and it is much cheaper to
    deploy than hardware attestation.
  \item[Tier 2, TEE path.] AGRP guarantees that specified
    code ran on hardware whose measurement the receiver can
    check against a vendor root. This is the claim
    proof-of-guardrail makes, and AGRP carries it as one
    backend rather than re-deriving it.
\end{description}

\paragraph{When a TEE is the right answer.}
Three conditions, and we recommend hardware attestation only
when at least one holds. First, the host operator is
untrusted: multi-tenant public infrastructure, or a
counterparty running the agent on infrastructure the
principal cannot audit. A mediator does not help here if the
mediator runs on the same untrusted host. Second, the
receiver's policy floor demands evidence of execution rather
than a claim, and no mediator is mutually acceptable. That
is the common case across a competitive boundary. Third, the
filter's own configuration is the thing at issue. A TEE
measurement covers the code and its configuration together,
while a mediator receipt covers only the observed
output. Absent all three, the mediator path buys the same
tier at a fraction of the operational cost, and we say so
rather than recommending hardware by default.

\paragraph{What a TEE does not buy, in AGRP's own terms.}
The proof-of-guardrail authors are explicit that their
scheme does not prevent a developer from deliberately
weakening the guardrail it then faithfully proves it
ran~\cite{proof_of_guardrail}. AGRP inherits that limit
exactly, and it is the reason the envelope carries
\texttt{filter\_measurement} against a receiver-side
allow-list rather than treating any attested filter as
acceptable. Attestation answers ``did the named filter
run?''. It cannot answer ``is that filter adequate?''. The
allow-list is where the second question is decided, and
\S\ref{sec:proposals:agrp:open} records that the governance
of that list is unsolved.

\paragraph{Normalizing receipts across vendors.}
The practical obstacle to a receiver accepting Tier~2
evidence is that the evidence is not interchangeable. An
Intel SGX quote, an AMD SEV-SNP report, an AWS Nitro
attestation document, and a mediator receipt in the Pipelock
format~\cite{pipelock} are four different objects. They
differ in encoding, in signature
algorithm, in what the trust root is, and in which fields
carry the measurement. A receiver that hard-codes one of
them is coupled to a vendor by its ingress policy. AGRP
therefore treats the vendor-specific artifact as an opaque
blob inside a thin normalizing wrapper, and specifies the
wrapper rather than the artifact. The wrapper carries five
fields. The first is the evidence kind, either
\texttt{tee\_quote} or \texttt{mediator\_receipt}. The second
is a vendor and format identifier, which names the
verification profile to apply. The third is the trust-root
identifier, which the receiver must already hold out of band.
The fourth is the measurement value, lifted into a common
encoding so that allow-list comparison never requires parsing
the vendor artifact. The fifth is the opaque artifact itself,
for full verification.

Two rules make the wrapper load-bearing rather than
decorative. First, a receiver MUST reject evidence whose
declared measurement disagrees with the measurement inside
the artifact. That prevents a sender from presenting an
allow-listed measurement over a different attestation.
Second, a receiver that holds no verification profile for the
declared vendor MUST treat the envelope as Tier~1, not as
Tier~2 with a warning. An unverifiable quote is exactly as
strong as a self-signed claim. This normalization is
specified here and unimplemented. The collaboration branch
of \S\ref{sec:proposals:acsp:pacms} verifies sender
signatures and measurement allow-lists, but it implements no
concrete vendor profile. No cross-vendor
interoperability result is therefore reported. Curating the profile
registry is the same governance problem as the filter
allow-list, and belongs in the same venue.

\subsubsection{Threat-model coverage against the 12-class catalogue}
\label{sec:proposals:agrp:threat}

The 12-threat catalogue of \S\ref{sec:rubric:axis8} and
\S\ref{sec:security:model} is the natural benchmark for
AGRP's spec-level coverage. We report AGRP's
addressed-count informally because the proposals are
not scored against the rubric
(\S\ref{sec:proposals}'s opening note).

AGRP \emph{addresses}: agent impersonation (the OWA-bound
signature in the attestation), message replay (the
envelope carries a sender-generated nonce in the
\texttt{filter\_signature}'s scope), message tampering
(the envelope is signed end-to-end), capability
escalation (the \texttt{classification.applied\_labels}
constrains what the receiver accepts), delegation token
theft (token-bearing fields are flagged
\texttt{secret:blocked} and rejected at egress),
registry poisoning (the OWA trust chain authenticates
the resolution path), audit-log tampering (the
redaction log is sender-signed, with independent Tier~2
or log-anchored evidence when receiver policy requires it), and
cross-agent privilege confusion (the
\texttt{boundary.kind} field is verified by the
receiver).

AGRP \emph{partially addresses}: prompt injection
(content-classification labels can mark a message as
adversarial-pattern-detected but the filter must run
before the LLM-receiver processes the content), and
side-channel privacy (the envelope can encrypt
metadata fields but the protocol does not specify a
metadata-privacy primitive normatively at the cutoff).

AGRP \emph{does not address}: side-channel privacy
specifics beyond metadata-field encryption,
agentcard supply-chain (which is an OWA concern, not an
AGRP one), denial-of-service (which is a transport
concern).

The total is 8 addressed, 2 partial, 2 not-addressed,
on a 12-class catalogue. The threat-coverage profile
approaches ACNBP's 10-of-12 (the highest in the
surveyed set) and substantially exceeds the 1-of-12 that
A2A, MCP, and ACP each score. The profile is achieved not by
re-specifying the protocol but by adding the envelope
extension.

\subsubsection{Open questions and limitations}
\label{sec:proposals:agrp:open}

\paragraph{Q1: Filter key management.}
The \texttt{filter\_signature} requires a key tied to
the agent's OWA URI. The key management discipline
covers key issuance, rotation, and revocation. It is not
specified in the AGRP envelope and is deferred to a future
revision. The natural choice is a binding to the
organization's DNSSEC-anchored trust chain through the
OWA \texttt{<org>} segment, but this couples AGRP to
DNSSEC deployment which not all organizations have.

\paragraph{Q2: Filter approval and allow-list governance.}
The receiver's allow-list of approved filters is the
locus of policy decision-making. An organization with
a single approved filter has a simple allow-list; a
federation of organizations needs a federated
allow-list. AGRP does not specify the allow-list
distribution protocol; the AAIF working group is a
plausible venue for that work in a future revision.

\paragraph{Q3: Backwards compatibility with non-AGRP A2A clients.}
An A2A client that does not speak AGRP receives a
message wrapped in the AGRP envelope and interprets
the envelope as an extension field it can ignore. The
\texttt{payload} field carries the unwrapped A2A
message in the format the client already understands.
The cost of ignoring AGRP is that the client doesn't
benefit from the egress-filtering guarantees; the
cost of \emph{not} ignoring is non-compliance with
the receiver's ingress policy. The AAIF interoperability
working group is the natural venue for specifying when
AGRP becomes a mandatory part of A2A rather than an
optional extension.

\paragraph{Q4: Filter misclassification, in both
directions.}
Every guarantee above is conditional on the filter's
classifier being right, and no content classifier is right.
This is the failure mode with the largest practical
consequence, so we name both directions and say which one
the design privileges.

A \emph{false negative} is a sensitive span the filter does
not recognise. The span crosses the boundary, and the
envelope attests that a filter ran, which is true. The
receiver, and any auditor reading the log later, sees a
clean attestation over a leaking message. This is strictly
worse than no envelope at all in one specific respect: an
unattested message invites suspicion, and an attested one
invites trust it has not earned. The design's answer is
partial and worth stating as partial. The envelope records
which filter and which version produced the log. A span
class later found to be missed can therefore be traced to
every message that filter version signed. That
converts a silent leak into a bounded, enumerable incident
after the fact. It does not prevent the leak, and the
reduced-log discipline of
\S\ref{sec:proposals:acsp:open} narrows what an
auditor can reconstruct across an organizational boundary.

A \emph{false positive} is a benign span the filter redacts.
The cost is task utility, and it is invisible in exactly the
way the false negative is not. The receiver gets a
well-formed message with a hole in it. The attestation is
valid, and nothing indicates that the redaction was wrong.
Under fail-closed semantics the pressure this creates is
the dangerous part. An operator whose agents fail often
enough will loosen the filter. A loosened filter that
still attests correctly produces exactly the false-negative
case above. The two failure modes are therefore coupled
through operator behaviour, and an evaluation that measures
only one of them will mis-state the risk.

AGRP privileges the false positive deliberately: the
envelope is fail-closed, and a span the filter cannot
classify is redacted rather than passed. That is the right
default for the regulated boundaries the proposal targets.
It is the wrong default for latency-sensitive or
utility-sensitive traffic. That is why the tier semantics
let \texttt{same\_workspace} traffic run at Tier~1 with a
weaker filter. What the specification does not do is
quantify either rate. Doing so requires a labelled corpus
of inter-agent messages with ground-truth sensitive spans,
and no such corpus exists at the cutoff. Reporting a
detection rate from an unlabelled corpus would be worse than
reporting none. Building that corpus is a prerequisite for
the benchmarking item P11 of
\S\ref{sec:openproblems:benchmarking}, not a detail of it.

\paragraph{Q5: Performance overhead.}
The envelope adds a per-message overhead: signature
computation, fingerprint hashing, and ingress
verification. The three backend choices
(\S\ref{sec:proposals:agrp:composition}) span three
orders of magnitude in expected cost: an in-process
filter adds roughly the cost of a single signature plus
the hashes of redacted fields; a proxy-style mediator
backend (Pipelock~\cite{pipelock} or equivalent) adds a
network round-trip on top of the in-process work; a TEE
backend (SAGA~\cite{saga2026} or equivalent) adds the
vendor-specific attestation overhead, which is typically
the largest of the three. The envelope itself does not
specify a target latency budget; AGRP is at
specification-sketch stage and no benchmarked
implementation exists at the cutoff. Establishing
quantitative overhead figures is deferred to the
benchmarking item P11 of
\S\ref{sec:openproblems:benchmarking}, which is the
natural venue for both AGRP and the broader
rubric-stable benchmark the field needs. One data point from
the adjacent literature bounds the hardware path. The
proof-of-guardrail authors report measured latency and
deployment cost for running an agent and its guardrail
inside a TEE~\cite{proof_of_guardrail}. That figure is not
AGRP's, since it measures a different envelope over a
different workload. It is still the closest published
estimate of what the Tier~2 hardware path costs.

\subsection{ACSP Context Selection Envelope (Contribution C7)}
\label{sec:proposals:acsp}

\subsubsection{Problem statement}
\label{sec:proposals:acsp:problem}

\S\ref{sec:security:privacy} documents that, within the
coordinator-worker configurations evaluated by AgentLeak,
conversation-history forwarding is a recurrent
inter-agent privacy leakage vector: the orchestrator hands
the remote agent prior context to spare it clarifying
questions, and the measured cost is that
\textbf{68.8\%} of inter-agent messages carry
unauthorised sensitive content under the benchmark's
ground-truth policy, with output-only audits
missing \textbf{41.7\%} of those violations
(AgentLeak~\cite{agentleak}). The contextual-integrity
framing of \S\ref{sec:security:privacy} isolates the root
cause: the message boundary is no longer the trust
boundary. The protocol implication
\S\ref{sec:security:privacy} leaves open is whether the
contextual-integrity primitive belongs at the protocol
layer or at the orchestrator's policy layer.

AGRP (Contribution C6, \S\ref{sec:proposals:agrp})
addresses one half of the gap by redacting spans within
forwarded content and attesting that the egress filter
ran. Span redaction, however, operates on what has
\emph{already been chosen} to forward; it cannot prevent
a sender from forwarding an irrelevant item, including an
item whose sensitive semantics are not captured by the
configured span detector. The remaining gap is the
\emph{selection} step that decides which items from the
sender's context pool are appropriate to cross the
boundary at all. Common implementations use recency
truncation (drop oldest items first) or top-k cosine
retrieval. Recency is not query-anchored; neither baseline
jointly models token cost, coverage, boundary policy, and
disclosure risk; and neither leaves a wire-level decision
record the receiver can audit.

ACSP (Contribution C7) is the protocol-layer answer to
the selection half of the contextual-integrity question.
The envelope carries the selection policy the sender
claims to have applied, the declared candidate log and
each disposition, and an attestation chain. A receiver can
authenticate the claim and check observable invariants
(cardinality, hard-gate decisions, mandatory conflicts,
and budgets). The envelope alone cannot prove that the
sender disclosed its complete pre-selection pool, executed
the claimed objective, or found an optimal set. The
reference implementation
(\S\ref{sec:proposals:acsp:pacms}) is the PACMS context
engine~\cite{pacms_cikm_demo}.

\subsubsection{Related work and differentiation}
\label{sec:proposals:acsp:related}

Table~\ref{tab:acsp-differentiation} compares ACSP
against eight prior systems and benchmarks in the
context-selection space.

\begin{table*}[t]
  \centering
  \caption{ACSP differentiation matrix. Among the eight
    comparators reviewed here, none is documented as
    simultaneously providing all four of
    query-anchored selection, budget-bounded selection
    with a declared solver, a wire-level
    audit log on the selection, and a boundary-keyed
    selection policy derived normatively from the
    sender-recipient naming.}
  \label{tab:acsp-differentiation}
  \footnotesize
  \resizebox{\textwidth}{!}{%
  \begin{tabular}{p{2.7cm}p{2.6cm}ccccc p{2.6cm}}
    \toprule
    Prior work & Mechanism & Query-anchored? & Budget-bounded solver? & Wire-level audit log? & Boundary-keyed? & Guarantee status & ACSP improvement \\
    \midrule
    Recency / last-k truncation
      & drop-oldest
      & no & no & no & no & n/a
      & topic-blind baseline; collapses under paraphrase redundancy \\
    Top-k cosine retrieval
      & pairwise similarity
      & yes & no & no & no & n/a
      & adds budget-bounded coverage and diversification \\
    MMR~\cite{mmr_carbonell}
      & pairwise penalty
      & yes & no & no & no & heuristic
      & coverage objective; explicit solver and wire-level envelope \\
    LangChain LC-MMR
      & production MMR
      & yes & no & no & no & heuristic
      & same gap as MMR; LangChain has no wire envelope \\
    BGE / cross-encoder rerankers
      & learned scoring
      & yes & no & no & no & learned
      & no budget primitive, no audit, no attestation \\
    RAG with budget caps
      & retrieve and truncate
      & partial & partial & no & no & n/a
      & principled selection inside the budget instead of cap-truncation \\
    Pipelock~\cite{pipelock}
      & egress proxy, span redaction
      & no (spans) & no & yes (receipt) & no & n/a
      & operates on items, not spans; complementary, not competing \\
    NeMo Guardrails
      & programmable rails
      & no & no & no & no & n/a
      & wire-level instead of orchestrator-level \\
    \bottomrule
  \end{tabular}%
  }
\end{table*}

Within the eight comparators and sources reviewed for the
matrix, the four properties of ACSP, namely
query-anchored selection, budget-bounded selection with a
declared solver, a wire-level audit log,
and a boundary-keyed policy derived from OWA naming,
appear individually but are not documented jointly. This
is a bounded comparison, not a claim that no other system
or later version has the combination. The four-property
argument is intentionally narrower than the
five-property argument used for OWA
(\S\ref{sec:proposals:owa}) and the seven-property
argument used for AGRP
(\S\ref{sec:proposals:agrp:related}); the co-design
property is deferred to \S\ref{sec:proposals:codesign} so
it is not double-counted.

Three cells of the table carry the argument most heavily.

\paragraph{MMR is the closest precedent and the sharpest contrast.}
The Maximal Marginal Relevance
algorithm~\cite{mmr_carbonell} is the closest precedent
in this comparison and is widely used in
retrieval-augmented systems. The contrast is
sharp on three axes. First, MMR's objective is
\emph{pairwise}: it maximises a tradeoff between query
relevance and pairwise dissimilarity from the
already-selected set. ACSP's reference objective is
\emph{facility-location coverage}, which sums over all
candidate items. The objective is monotone submodular, but
a constant-factor claim for a token-knapsack implementation
depends on the exact algorithm and assumptions; the
wire-level record does not establish that proof. Second, MMR's
\(\lambda\) parameter is heuristic; ACSP's budget is a
token knapsack with a principled solver. Third, MMR is a
library function; ACSP is a wire envelope. The receiver
does not see what MMR considered or rejected; the
receiver sees the ACSP selection log as a matter of
normative specification.

\paragraph{Pipelock is complementary, not competing.}
Pipelock~\cite{pipelock} is the deployment-layer state of
the art for content redaction at egress, and the same
prior work AGRP positions as complementary
(\S\ref{sec:proposals:agrp:composition}). Pipelock and
ACSP operate on \emph{different units}: ACSP selects
\emph{items} (candidate turns, memories, tool outputs)
that may enter the forwarded context; Pipelock and AGRP
redact \emph{spans} (matched patterns within selected
text). The two are sequential, not substitutive: ACSP
narrows the candidate set; AGRP / Pipelock cleans the
spans that remain in the items ACSP kept. Together with
OWA's boundary detection, the three compose into the
triplet of \S\ref{sec:proposals:codesign}.

\paragraph{NeMo Guardrails sits at the wrong layer.}
NeMo Guardrails is the most-developed
orchestrator-layer programmable-rails system. The
orchestrator declares allow/deny rules against the
conversational state and NeMo enforces them inside the
orchestrator process. \S\ref{sec:security:privacy} calls
out the orchestrator-versus-protocol dispatch
explicitly: NeMo represents the orchestrator side of
that dispatch. ACSP is the protocol-side counterpart.
The two are not in conflict. A NeMo-equipped
orchestrator can emit ACSP envelopes as part of its
outgoing-message rule. But the question of which
layer the primitive \emph{normatively} belongs at is
what \S\ref{sec:security:privacy} leaves open, and ACSP
is the survey's protocol-layer answer.

\subsubsection{The ACSP envelope: schema and semantics}
\label{sec:proposals:acsp:envelope}

The ACSP envelope is a JSON-shaped block sibling to AGRP
inside the A2A message wrapper. The sender produces ACSP
before running AGRP over the kept items. The receiver
verifies AGRP first and then ACSP, as specified in
\S\ref{sec:proposals:acsp:composition}; it processes no
payload until both required checks succeed.

\begin{verbatim}
{
  "acsp_version": "1.0",

  "sender":    "agent+owa://acme.com/w/finance/c/billing/a/agent-01",
  "recipient": "agent+owa://acme.com/w/support/c/chat/a/agent-02",

  "boundary": {
    "kind": "cross_workspace"
  },

  "selection_policy": {
    "policy_id":           "acme/finance.context.v2",
    "objective":           "submodular_facility_location",
    "budget_tokens":       3800,
    "risk_budget":         1.25,
    "risk_weight":         0.50,
    "mandatory_set_kind":  "plan_edges",
    "hard_gate":           "credential-egress-v1"
  },

  "selection_log": {
    "pool_size":      2,
    "selected":       1,
    "dropped":        0,
    "hard_blocked":   1,
    "dropped_tokens": 38,
    "commitment_scheme": "HMAC-SHA-256",
    "commitment_status": "present",
    "commitment_key_id": "acme-context-key-2026q3",
    "commitment_nonce": "4f7c...32-hex-chars...",
    "declared_pool_commitment": "hmac-sha256:a123...",
    "candidates": [
      { "index": 0, "candidate_ref": "candidate-0000",
        "kind": "tool_output", "classification": "internal",
        "tokens": 142, "mandatory": false,
        "selected": true, "decision": "selected",
        "risk_score": 0.08,
        "commitment": "hmac-sha256:9d3f..." },
      { "index": 1, "candidate_ref": "candidate-0001",
        "kind": "user_message", "classification": "restricted",
        "tokens": 38, "mandatory": false,
        "selected": false, "decision": "hard_blocked",
        "risk_score": 1.00,
        "commitment": "hmac-sha256:c41a..." }
    ]
  },

  "attestation": {
    "selector_id":          "pacms-v0.1.0",
    "selector_measurement": "sha256:6ec4...64-hex-chars...",
    "signing_key_id":       "owa-key-2026q3",
    "assurance_tier":       1,
    "evidence_status":      "self_attested_claim",
    "selector_signature":   "ed25519:base64url...",
    "mediator_receipt":     "... optional ...",
    "tee_quote":            "... optional ..."
  },

  "payload": { /* A2A message carrying only the selected items */ }
}
\end{verbatim}

The fields carry the following semantics. The
\texttt{sender} and \texttt{recipient} fields are OWA
URIs (Contribution C5). The \texttt{boundary.kind} field
is derived by prefix comparison of the sender and
recipient URIs, identical to AGRP's procedure
(\S\ref{sec:proposals:agrp:boundary}); the same boundary
detection routine serves both envelopes.

The \texttt{selection\_policy} block declares the
selection contract the sender claims to have honoured.
The \texttt{objective} field is an enum of admissible
optimisation objectives; the protocol specifies the enum,
not the implementation. At the 2026-Q1 cutoff the enum
contains \texttt{submodular\_\allowbreak{}facility\_\allowbreak{}location} (the
PACMS reference; coverage maximisation under a token
knapsack),
\texttt{mmr\_lambda} (MMR with declared
\(\lambda\); heuristic, no approximation guarantee), and
\texttt{topk\_cosine} (no diversification). The
\texttt{budget\_tokens} field is the token knapsack
budget. The \texttt{mandatory\_set\_kind} field declares
how the sender derived the items the selector was not
permitted to drop. Its values are \texttt{plan\_edges}
(items the sender's plan declared task-critical),
\texttt{system\_prompt}, \texttt{user\_pinned} (items the
user marked must-include), and \texttt{none}.

The \texttt{selection\_log} is the operationally
consequential field. The \texttt{pool\_size} declares the
candidate count the selector saw; \texttt{selected},
\texttt{dropped}, and \texttt{hard\_blocked} declare the
three decision counts; \texttt{dropped\_tokens} is the gross
token volume the selector did not select. The
\texttt{candidates} array records each candidate's index,
kind (\texttt{user\_message}, \texttt{assistant\_message},
\texttt{memory}, \texttt{tool\_output}), classification,
token count, mandatory flag, decision
(\texttt{selected}, \texttt{dropped},
or \texttt{hard\_blocked}), optional utility score, risk
score, and a keyed HMAC commitment. A bare content hash is
insufficient: low-entropy candidate text is vulnerable to
dictionary recovery. The \texttt{commitment\_key\_id}
selects a key provisioned to an authorised verifier by an
out-of-band mechanism; ACSP~1.0 does not define key
negotiation. The \texttt{commitment\_nonce} is fresh for
each envelope and is included in the HMAC input to prevent
replay of an otherwise valid commitment. If no key is
configured, the implementation records
\texttt{key\_not\_configured} and omits the commitment
rather than silently emitting an unkeyed hash. The
per-candidate commitment binds candidate content, index,
and kind. The declared-pool commitment additionally binds
those commitments, candidate decision metadata, and policy
context. Together they make those declared fields
tamper-evident to a verifier holding the key; they do not
prove that the selector optimised the declared objective or
that the sender disclosed its complete pre-selection pool.
AgentLeak's empirical
finding that output-only audits miss 41.7\% of violations~\cite{agentleak}
motivates evidence beyond final output; it does not by
itself establish that the absence of a selection envelope
caused those missed violations.

The \texttt{attestation} block records the sender's signed
selection claim and optional independent evidence that the
selector ran. The \texttt{selector\_id} and
\texttt{selector\_measurement} identify the claimed
selector binary; the \texttt{selector\_signature} is
signed by the sender's OWA-bound key. That signature alone
is Tier~1 self-attestation: it authenticates the claim but
does not prove execution. The optional
\texttt{mediator\_receipt} and \texttt{tee\_quote} fields
reuse AGRP's Tier~2 independence and verification rules
(\S\ref{sec:proposals:agrp:tiers}). This parallel is
deliberate: a receiver that verifies AGRP envelopes can
reuse the same trust roots and cryptographic infrastructure
for ACSP envelopes.

The receiver's verification procedure on an inbound ACSP
envelope, \emph{fail-closed} on any mismatch, comprises
eight checks:

\begin{enumerate}
  \item Verify the \texttt{selector\_signature} against
    the sender's OWA-bound public key. If invalid, fail
    closed.
  \item Check the \texttt{selector\_measurement} against
    the receiver's allow-list of approved selectors for
    the boundary kind. If the selector is not on the
    allow-list, fail closed.
  \item Determine and verify the attestation tier using
    AGRP's independence rules. If receiver policy requires
    Tier~2, validate the mediator receipt or TEE quote
    against a trust root provisioned out of band; a
    sender-only signature MUST NOT satisfy that policy.
  \item Cross-check the cardinality of the
    \texttt{selection\_log}: the count of
    \texttt{selected}, \texttt{dropped}, and
    \texttt{hard\_blocked} candidates must equal
    \texttt{pool\_size}. If not, fail closed.
  \item Validate the budget: the sum of \texttt{tokens}
    over the \texttt{selected} candidates must not exceed
    \texttt{budget\_tokens}. If exceeded, fail closed.
  \item Apply the receiver's ingress policy to the
    \texttt{mandatory\_set\_kind} field and per-candidate
    \texttt{mandatory} flags. Every declared mandatory
    candidate must have decision \texttt{selected}; a
    mandatory candidate marked \texttt{dropped} or
    \texttt{hard\_blocked} is a
    fail-closed conflict. Proving that the sender omitted no
    mandatory candidate requires an independently anchored
    mandatory-set manifest and is outside envelope-only
    verification.
  \item Verify keyed candidate commitments when the
    boundary policy requires them. A receiver that lacks
    the verification key records the check as unavailable;
    it MUST NOT treat an unkeyed hash as equivalent.
  \item Check per-objective invariants that are observable
    from the envelope (hard-gate exclusions, token and risk
    budgets, and mandatory-set conflicts). Recomputing
    optimality or an approximation ratio requires the full
    candidate pool, scoring model, and solver configuration
    and is outside envelope-only verification.
\end{enumerate}

Fail-closed is normative: an ACSP-aware receiver MUST
refuse to process the payload if any of the eight
verification steps fail. A receiver that downgrades to
``process anyway with a warning'' breaks the protocol's
guarantees, and SHOULD be flagged non-compliant by an
ACSP conformance test suite.

\subsubsection{Security-gated selection and four-channel evaluation}
\label{sec:proposals:acsp:security-eval}

The ACSP envelope makes the selection decision visible,
but visibility alone does not make the decision safe. We
therefore add a security-gating profile that separates
\emph{non-negotiable constraints} from \emph{graded risk}.
This distinction avoids two opposite errors: allowing a
high-utility credential because its relevance score is
large, and blocking every benign discussion that contains
the word ``password.''

\paragraph{Hard and soft gating.}
Let $V$ be the candidate pool, $c_i$ the token cost of
candidate $i$, $F_q(S)$ the query-conditioned PACMS
coverage utility, $H_b(i)$ an allow/deny hard-policy
decision for boundary $b$, and $r_b(i) \in [0,1]$ a soft
risk score. The security-aware selection problem is

\begin{equation}
  \max_{S \subseteq V}
    F_q(S) - \lambda_r \sum_{i \in S} r_b(i)
  \quad \text{s.t.} \quad
    \sum_{i \in S} c_i \leq B,\quad
    \sum_{i \in S} r_b(i) \leq R,\quad
    H_b(i)=\mathrm{allow}\;\forall i\in S.
  \label{eq:acsp-security-selection}
\end{equation}

In the envelope, the \texttt{risk\_budget} field in
\texttt{selection\_policy} instantiates $R$, while its
\texttt{risk\_weight} field instantiates $\lambda_r$.

The prompt has its own hard decision $H_b(q)$. If
$H_b(q)=\mathrm{deny}$, the request is rejected before selection. A
candidate for which $H_b(i)=\mathrm{deny}$ is excluded even if the
orchestrator marked it mandatory: mandatory status cannot
override a confidentiality policy. The current PACMS
prototype realises Equation~\ref{eq:acsp-security-selection}
as a risk penalty on marginal utility plus independent
token and cumulative-risk feasibility checks. We describe
this as a solver implementation, not as a new approximation
guarantee.

The deterministic first tier hard-blocks explicit
credential-exfiltration intent, assigned credential values,
and private-key material. Credential topics and email
addresses contribute soft risk. A benign request about key
rotation is therefore scored but not rejected unless a
deployment explicitly chooses a keyword-only policy. These
rules are auditable and reproducible, but intentionally
incomplete: paraphrase, multilingual intent, encoded
secrets, and domain-specific privacy semantics require
additional detectors and red-team evaluation. The policy
interface accepts their scores without coupling the
selector to one detector implementation.

\paragraph{Lifecycle stages and disclosure channels.}
Selection-time controls see the request and candidate pool;
they do not observe every later disclosure. The model and
operational runtime can disclose an item that was safely
selected, transform it into a newly sensitive value, or
leak it through a fallback. The evaluation therefore uses
a two-dimensional model. Lifecycle stages are input,
candidate-pool construction, selection, execution, and
output. AgentLeak channels describe where data is observed:
\texttt{user\_input} and \texttt{tool\_response} are
sources, while the four channels of
Table~\ref{tab:acsp-four-channel} are the core disclosure
frame.

\begin{table}[t]
  \centering
  \caption{Core disclosure channels for ACSP security
    evaluation. A channel
    marked not observed is missing evidence, not a clean
    result.}
  \label{tab:acsp-four-channel}
  \footnotesize
  \begin{tabular}{p{3.0cm}p{3.0cm}p{4.2cm}}
    \toprule
    Disclosure surface & AgentLeak channel & Question \\
    \midrule
    Agent handoff & \texttt{inter\_agent\_message} & Did a selected item reach another agent across the declared boundary? \\
    Persistent context & \texttt{shared\_memory} & Was a selected or generated sensitive item persisted for later agents or turns? \\
    Tool arguments & \texttt{tool\_call} & Did sensitive content enter an invocation payload or external action? \\
    User-facing answer & \texttt{final\_output} & Did selected or newly generated sensitive content reach the principal? \\
    \bottomrule
  \end{tabular}
\end{table}

The \texttt{log} and \texttt{generated\_file} channels are
additional operational surfaces when the runtime exposes
them. The selection API emits a content-free decision
record for the input and candidate pool, not an AgentLeak
channel verdict. A later trace join links each selected
\texttt{candidate\_ref} to zero or more disclosure-channel
findings. This status discipline prevents an unavailable
detector or absent trace from being counted as a
non-disclosure.

\paragraph{Benchmark protocol and metrics.}
The post-cutoff collaboration artefact includes a
deterministic feasibility harness with controlled canaries,
fixed embeddings, the real PACMS selector, and the real
AgentLeak trace analyser. It exercises the four core
disclosure channels plus the \texttt{log} and
\texttt{generated\_file} paths. Its deterministic
forced channel-injection simulator is deliberately not an
LLM and does not model runtime causality. It validates the measurement path with three
arms: no controls, hard gating only, and hard plus soft
gating. The generated report carries the evidence label
\texttt{synthetic\_harness\_validation} and contains no raw
canary values. It MUST NOT be interpreted as evidence of
PACMS security effectiveness, LLM utility, OpenClaw
fidelity, or a Nasiko deployment.

The empirical study will replace the forced channel-injection simulator with a
version-locked reader and will cross the security arms with
PACMS, top-k, last-k, MMR, and fixed-$N$ RAG. For each
scenario and seed it reports, separately:

\begin{enumerate}
  \item required-context recall and end-task utility;
  \item unauthorised-sensitive keep rate at selection;
  \item disclosure incidence and severity by observed
    channel;
  \item malicious-request block rate and benign-request
    false-positive rate;
  \item selected tokens, latency, and cumulative risk;
  \item fail-closed versus explicitly configured fail-open
    behavior under sidecar and logging faults.
\end{enumerate}

AgentRisk is not yet published as an independent method.
Its global, per-channel, and level-profile values may be
retained under an explicitly labelled exploratory field,
but they are not cited or used as primary evidence in this
study. Primary disclosure and utility metrics are defined
directly in the protocol above.

No single scalar combines these outcomes. In particular, a
privacy improvement does not compensate for unusable task
performance, and a utility improvement does not compensate
for a hard-policy violation. Effect estimates are paired by
scenario and seed with confidence intervals. A preregistered
capture-completeness rule marks incomplete observations
\texttt{not\_observed}, excludes them from the primary
estimator, and reports missingness plus conservative
sensitivity analyses. Captured selector or runtime failures
remain in the intention-to-test denominator.

\paragraph{Integration with the survey rubric.}
Axis~8 remains the count of threats addressed at the
\emph{specification} level. A PACMS deployment test must not
increase a protocol's Axis~8 count. The four-channel record
is instead an empirical companion to three catalogue items:
side-channel privacy leakage (Threat~6), foundation-model
prompt injection (Threat~10), and audit-trail tampering or
omission (Threat~11). Each benchmark record binds the threat
identifier, boundary, policy version, selector version,
experimental arm, and evidence status. This lets the
rubric-based comparison distinguish ``the specification
addresses the threat'' from ``this implementation resisted
these controlled cases.''

\paragraph{Implementation boundary.}
\begin{itemize}
  \item \textbf{Implemented.} The collaboration branch
    implements deterministic hard and soft assessments,
    token and risk constraints, MMR parity, nonce-bound
    content-free HMAC candidate and declared-pool
    commitments, a portable logger with explicit
    audit-failure behavior, a content-free OTLP selection
    lineage exporter, and the synthetic PACMS/AgentLeak
    harness. Its OpenClaw plugin is version-pinned,
    configuration-validated, timeout-bounded, buildable, and
    fail-closed by default. The branch also contains a Nasiko
    adapter on both production history-assembly paths; it is
    disabled when no PACMS URL is configured, uses the same
    explicit failure modes, ignores selector-returned text in
    favour of local messages, and emits content-free
    selection spans. A shared RFC~8785/Ed25519 profile signs
    complete ACSP and AGRP envelopes and verifies sender
    signatures, measurement allow-lists, declared tiers, and
    observable ACSP invariants. Its default is explicitly
    Tier~1 self-attestation.

  \item \textbf{Tested.} Unit, HTTP-contract, TypeScript,
    mocked lifecycle, Rust adapter, cryptographic, and
    synthetic integration checks exercise the stated paths.
    A real pinned OpenClaw Gateway smoke loads the plugin as
    the active context engine and passes authenticated
    health. It performs no model call and is therefore not
    an end-to-end fidelity or effectiveness result. Nasiko
    unit tests exercise configuration, HTTP contract,
    response validation, and fallback semantics; no deployed
    Nasiko experiment is reported.

  \item \textbf{Still gated.} A Tier~2 result requires a
    caller-supplied independent evidence verifier; the branch
    does not implement a concrete mediator-receipt or TEE
    vendor verifier, OWA key discovery, automatic sidecar
    envelope emission, or remote attestation. A
    version-locked model run, deployed Nasiko run, and fresh
    LongMemEval/security evaluation remain gates before any
    effectiveness claim.
\end{itemize}

\subsubsection{Composition with OWA and AGRP}
\label{sec:proposals:acsp:composition}

ACSP composes \emph{upstream} of AGRP. The processing
pipeline at the sender is:

\begin{enumerate}
  \item OWA's boundary detection
    (\S\ref{sec:proposals:agrp:boundary}) classifies the
    sender-recipient pair as \texttt{same\_workspace},
    \texttt{cross\_workspace}, or \texttt{cross\_org}.
  \item The sender's selection-policy table maps the
    boundary kind to an ACSP policy (objective, budget,
    mandatory-set kind).
  \item The selector runs over the candidate pool and
    emits the ACSP envelope's selection log.
  \item AGRP's egress filter runs over the items ACSP
    kept and emits the AGRP envelope's redaction log.
  \item Both envelopes are bound to the same OWA-rooted
    identity and travel with the A2A message.
\end{enumerate}

The receiver verifies in reverse order: AGRP first (was
the egress filter run on what arrived?), then ACSP (was
the selection that produced the arriving items
consistent with the declared policy?), then the payload
itself. The two envelopes are normatively independent:
an A2A message can carry AGRP without ACSP or vice
versa. But \S\ref{sec:proposals:codesign} argues that
the deployment posture that closes the
contextual-integrity loop is the one carrying both
together with OWA.

\subsubsection{Reference implementation: PACMS}
\label{sec:proposals:acsp:pacms}

The PACMS system~\cite{pacms_cikm_demo} is an operational context
selector that instantiates part of the contract ACSP
specifies. PACMS formulates context selection as
budget-constrained submodular maximisation of a
facility-location coverage objective, solved by a CELF
lazy-greedy procedure. The paper associates the objective
with the classical submodular guarantee; the present survey
does not independently prove that the specific
token-knapsack implementation satisfies all conditions
required for that bound. The mandatory-set abstraction
(\emph{plan edges}) maps directly to the
\texttt{mandatory\_set\_kind} field of the ACSP envelope's
selection policy.

PACMS is packaged as a FastAPI sidecar that exposes a
\texttt{POST /select} endpoint. The post-cutoff
collaboration branch adds risk inputs and a content-free
decision ledger, bringing the request and response closer
to ACSP's \texttt{selection\_policy} and
\texttt{selection\_log} fields. A separate module implements
the Tier~1 ACSP/AGRP signer and core receiver verifier, but
the \texttt{/select} endpoint does not yet emit a complete
signed ACSP wire envelope. Concrete Tier~2 evidence
verification, OWA key discovery and boundary derivation,
and receiver deployment remain unimplemented; PACMS
therefore remains a reference selector plus a partial
conformance profile, not a complete ACSP deployment.

PACMS reports a LongMemEval evaluation of retrieval and
downstream QA~\cite{pacms_cikm_demo}. During the
collaboration audit, the PDF, adjacent TeX source, and
runner were found to encode non-identical protocol details
(including strategy support, model configuration, and
redundancy injection). A canonical configuration and a
fresh reproduction are therefore required before this
survey reuses any numerical result. The security subsection
reports no LongMemEval or model result.

The audited OpenClaw prototype originally failed open when
the sidecar was unreachable. The collaboration branch now
defaults to \emph{fail-closed}: no historical context is
forwarded when selection fails. An operator can explicitly
choose an availability-first fail-open mode, which the
benchmark treats as a distinct experimental condition.
The ACSP envelope specification likewise requires
\emph{fail-closed at the receiver}
because a verifiable wire-layer primitive whose absence
is interpreted as ``no selection happened'' would amount
to a forwarding-by-default policy at exactly the layer
\S\ref{sec:security:privacy} identifies as the leakage
vector. Sender fallback and receiver verification remain
separate decisions, but neither is silently inferred.

\subsubsection{Open questions and limitations}
\label{sec:proposals:acsp:open}

Four open questions warrant explicit treatment.

\paragraph{Q1: Selector-binary approval governance.}
ACSP's receiver verification (step 2) requires the
receiver to maintain an allow-list of acceptable
selector binaries by \texttt{selector\_measurement}.
Within a single organisation this is tractable. Across
an inter-organisational boundary it requires either
(a)~an industry-curated allow-list maintained by the
AAIF or W3C working groups, or (b)~a federation pattern
in which each organisation publishes its acceptable
measurement set through a DNS-rooted manifest the
receiver can discover. Neither pattern is specified at
the cutoff, and the choice between them is the same
governance question OWA leaves open for inter-org trust
anchoring (\S\ref{sec:proposals:owa}).

\paragraph{Q2: Objective enum extensibility.}
The \texttt{objective} enum at the cutoff is small
(\texttt{submodular\_\allowbreak{}facility\_\allowbreak{}location},
\texttt{mmr\_lambda}, \texttt{topk\_cosine}). A new
admissible objective requires a spec revision; the
working-groups path (AAIF / W3C) is the natural venue.
The open question is whether the spec should additionally
reserve a prefixed namespace for organisation-specific
objectives
(\texttt{x-acme/learned-reranker-v3}) that a receiver
can opt to accept. The trade-off is between
extensibility at deployment time and the cross-vendor
verifiability the spec aims to provide.

\paragraph{Q3: The selection log discloses candidate-set
membership, which is a privacy cost of the design.}
The envelope's central claim is that a receiver can check a
selection it did not perform. That check requires the
receiver to learn something about items the sender chose not
to send, and this is a real disclosure, not an artefact of a
sloppy schema. It should be stated as a cost of the design
rather than left for a reader to notice.

Concretely, an ACSP envelope tells the receiver at least
four things about the withheld items. It says how many there
were (\texttt{pool\_size} minus \texttt{selected}). It says
what kind each was (\texttt{user\_message}, \texttt{memory},
\texttt{tool\_output}). It says what classification band the
sender assigned it, and how many tokens it occupied. It also
says which were \texttt{hard\_blocked} rather than merely
\texttt{dropped}. That is the most informative field of the
set, because a hard block announces that the sender held an
item its own policy marks as credential-bearing or regulated.
Three attacks follow. A receiver can \emph{probe}. By varying
its request and watching \texttt{pool\_size} and the
classification histogram move, it learns which topics cause
the sender's pool to contain restricted material. It can
\emph{fingerprint}. Token counts plus kinds plus ordering
form a distinguishing signature. A receiver that sees the
same sender repeatedly can then tell whether a particular
sensitive item is still in the pool from one session to the
next. And it can \emph{confirm}. Given an external guess
about a specific item, a matching token count and
classification raises confidence in the guess considerably,
even though the content itself never crossed.

None of this is fixed by the keyed commitments. Those
protect integrity of the declared record, and the leakage
here is in the metadata the record is supposed to declare.
The tension is structural: verifiability of an omission and
confidentiality of what was omitted pull in opposite
directions, and the envelope as specified resolves that
tension toward verifiability.

ACSP~1.0 therefore specifies the leakage as
boundary-conditional rather than pretending it away. The
full per-candidate array is permitted for
\texttt{same\_workspace} traffic. For
\texttt{cross\_workspace} and \texttt{cross\_org}
boundaries, the sender MAY emit a \emph{reduced log}. It
keeps only the fields the receiver's checks actually need:
the cardinality counts, the aggregate
selected-token total against the budget, the mandatory-set
conflict flag, and the declared-pool commitment. It omits
per-candidate kinds, classifications, and token counts.
Every check of \S\ref{sec:proposals:acsp:envelope} except
the per-candidate commitment verification survives that
reduction, because those checks are aggregate. The
per-candidate commitments remain available to an auditor
holding the key, out of band, which is the party that
actually needs them.

Three limits on this mitigation, stated plainly. The reduced
log still discloses cardinality, and cardinality alone
supports the probing attack in weakened form. We have not
quantified the residual leakage. Doing so needs a defined
adversary and a measured experiment, and no such experiment
is reported here. And a receiver that requires the full log
before it will trust a sender can force the disclosure by
policy, so the choice is not entirely the sender's. A
privacy-preserving proof of the aggregate counts would let
the receiver check cardinality without learning it. Whether
that is the right answer is open, and it is the most
interesting open question the envelope raises.

\paragraph{Q4: Cross-org embedding-model agreement.}
A receiver cannot independently recompute selection
scores if the sender's embedding model is private or
proprietary. The keyed-commitment discipline is the fallback:
the receiver verifies that the selection log is
internally consistent with the declared scores and that
mandatory items are present, but does not re-run the
selector. Whether this is sufficient for high-assurance
deployments (regulated finance, clinical multi-agent
systems) or whether the spec should additionally require
a \texttt{model\_attestation} field analogous to AGRP's
\texttt{filter\_measurement} remains open. The trade-off
is between independent verifiability and IP protection
of the sender's embedding model.

ACSP, like OWA and AGRP, is at specification-sketch
stage at the 2026-Q1 cutoff: the PACMS reference
implementation exists, but no inter-organisational
deployment has been benchmarked end-to-end. The
benchmarking item P11 of
\S\ref{sec:openproblems:benchmarking} is the natural
venue for both ACSP and the broader triplet of
\S\ref{sec:proposals:codesign}.

\subsection{Why C5, C6, and C7 are co-designed}
\label{sec:proposals:codesign}

The closing observation of \S\ref{sec:proposals} is that
the three proposals are co-designed. They specify a
single load-bearing operation, \emph{naming, selection,
attestation}, at three layers of one loop. And no single
piece preserves the invariant the triplet
preserves.

OWA without ACSP and AGRP gives an inter-agent naming
scheme with a workspace tier but no protocol-layer
enforcement of the policy boundaries the workspace
identifies. The naming improvement is real (the
workspace tier is absent from \texttt{agent://}, ANS,
A2A, and DIDs) but the practical adoption case is weaker.
An organisation that wants per-workspace policy must
implement both \emph{which items cross} and \emph{which
spans are cleaned} at the deployment layer. That is
where the contemporary protocols already leave them.

ACSP without OWA and AGRP gives a query-anchored,
budget-bounded selection envelope with a verifiable
selection log, but no normative way to derive what
policy the selection should apply (the boundary kind
that keys the policy is OWA-derived) and no companion
span-level cleaning of the items it kept. The envelope
would either specify its own boundary notion (duplicating
OWA) or accept span-level leakage in the items it
forwarded (re-opening the gap
\S\ref{sec:security:privacy} documents).

AGRP without OWA and ACSP gives a span-level redaction
envelope with attestation, but no way to derive the
boundary kind (OWA's role) and no way to minimise the
items forwarded in the first place (ACSP's role). Span
redaction handles the \emph{a selected item carries a
recognised PII span} case but not an item whose sensitive
semantics fall outside the configured detector, nor a
later disclosure through an internal operational channel.
Those channels motivate the four-channel evaluation of
\S\ref{sec:proposals:acsp:security-eval}.

Together, the three proposals close a single loop:
\textbf{OWA names the boundary; ACSP selects which items
cross it; AGRP redacts the spans within those items}.
Each step is normatively specified, each step is
receiver-verifiable, and the chain composes into the
invariant that \S\ref{sec:security:privacy} calls for and
that no contemporary protocol provides:
\emph{contextual integrity at the wire layer, end-to-end
and receiver-verifiable}. The triplet is the
protocol-layer answer to the question
\S\ref{sec:security:privacy} leaves open. That question is
whether the contextual-integrity primitive belongs at the
protocol layer or at the orchestrator's policy layer. The
triplet answers it without
pre-closing the orchestrator-side alternatives the
open-problems agenda catalogues
(\S\ref{sec:openproblems:regulation}, P5).

The adoption case for the triplet is the case
\S\ref{sec:adoption:healthcare} to \ref{sec:adoption:finance}
identifies as the principal cross-organisational
deployment blocker: organisations in regulated verticals
need per-workspace contextual-integrity policies they
can verify at the receiver's side. OWA gives them the
name; ACSP gives them the selection contract; AGRP
gives them the span-level verification. No pair alone is
sufficient; together they are.

The survey's broader argument
(\S\ref{sec:governance:survival}) is that the AAIF
working groups are the institutional venue most likely
to elevate the contemporary protocols toward closing the
open problems of \S\ref{sec:openproblems}. The triplet
of proposals here is the survey's contribution to that
venue: three specification-sketches the working groups
can adopt, extend, or reject, derived from the gaps the
rubric identifies and the deployment evidence
\S\ref{sec:adoption} catalogues.
\setlength{\emergencystretch}{0pt}
